\documentclass[11pt,oneside]{book}

% ---------- Encoding & fonts ----------
\usepackage[utf8]{inputenc}
\usepackage[T1]{fontenc}
\usepackage{textcomp}

% ---------- Math ----------
\usepackage{amsmath}
\usepackage{amssymb}
\usepackage{amsthm}

% ---------- Page layout ----------
\usepackage[margin=1.25in]{geometry}
\usepackage{setspace}
\onehalfspacing

% ---------- Quotes / typography ----------
\usepackage{csquotes}
\usepackage[protrusion=true,expansion=false]{microtype}

% ---------- Links (load late) ----------
\usepackage{hyperref}
\hypersetup{
  colorlinks=true,
  linkcolor=black,
  citecolor=black,
  urlcolor=black,
  pdftitle={The Crinkle-Cut Supercube},
  pdfauthor={Flyxion}
}

% ---------- Section numbering depth ----------
\setcounter{secnumdepth}{2}
\setcounter{tocdepth}{2}

% ---------- Quote environment styling (used for the Crinkle-Cut Principle, etc.) ----------
\usepackage{etoolbox}
\AtBeginEnvironment{quote}{\itshape}

% ---------- Pandoc compatibility macros ----------
\providecommand{\tightlist}{%
  \setlength{\itemsep}{0pt}\setlength{\parskip}{0pt}}

\title{\textbf{The Crinkle-Cut Supercube}\\[0.5em]
\large Constraint-Directed Infolding, Hidden Surface Area,\\
and the Geometry of Emergent Boundaries}
\author{Flyxion}
\date{}


\begin{document}
\pagestyle{plain}

\frontmatter

\maketitle

\hypertarget{abstract}{%
\chapter*{Abstract}
\addcontentsline{toc}{chapter}{Abstract}\label{abstract}}

Ordinary origami specifies a transformation and treats the resulting
boundary as its consequence: a sequence of folds, or the crease pattern
that records their result, is given in advance, and the finished shape
follows from executing it. This monograph describes and formalizes a
different construction, in which a small number of relations the
finished boundary must satisfy are specified instead, and the
transformation that satisfies them is discovered rather than prescribed.

The central example is a physical construction, the Crinkle-Cut
Supercube, in which three non-contiguous squares were marked on a single
large sheet, each annotated with the edge and vertex roles it was to
occupy in a finished cube corner. No crease pattern connected the three
squares in advance. The sheet was folded, pleated, and where necessary
crumpled until the three marked regions achieved their declared exterior
relations, with the remainder of the sheet, considerably larger in area
than the three faces it supports, folded into the object's interior
rather than trimmed away.

The monograph argues that existing origami notation, built around
sequences of operations or complete crease patterns, cannot describe a
specification of this kind, and introduces a fourth descriptive level,
the constraint level, together with a formal vocabulary of marked
regions, local frames, boundary grammars, and a constrained search
operator, sufficient to state such a specification precisely. This
formalism is then used to establish the monograph's central theoretical
claim: that a finished boundary satisfying a declared specification does
not, and structurally cannot, disclose which of the many interior
histories capable of satisfying that same specification actually
produced it. Exterior equivalence, in the sense the monograph makes
precise, does not imply structural, crease, or procedural equivalence
among the realizations that share it.

The argument is extended to the embodied character of the search process
itself, to a carefully bounded structural analogy with the folding of
the cerebral cortex, and to a formally hedged principle describing the
conditions under which spatially separated, locally framed regions can
be brought into a shared exterior configuration without cutting the
sheet that connects them. Throughout, the monograph distinguishes what a
single constructive demonstration actually establishes from what remains
conjectural, and treats the finished, visibly irregular object not as an
imperfect execution of a clean design but as a physical certificate that
a sparse specification, realized by a much larger and largely
unrecoverable hidden process, was in fact satisfiable.

\tableofcontents

\mainmatter

\hypertarget{introduction-the-boundary-before-the-fold}{%
\chapter*{Introduction: The Boundary Before the Fold}
\addcontentsline{toc}{chapter}{Introduction: The Boundary Before the Fold}\label{introduction-the-boundary-before-the-fold}}

Ordinary origami begins with a transformation and ends with a shape. A
sequence of folds, or a crease pattern that records their result, is
specified in advance, and the folder's task is to execute it faithfully.
The final boundary of the object is a consequence of the folds, not
their premise.

Constraint-directed infolding inverts this order. Instead of specifying
the transformation and observing what boundary results, the constructor
specifies a small number of relations the final boundary must satisfy
and permits the transformation itself to be found rather than
prescribed. Selected regions of a sheet are marked and assigned roles:
this edge is to become a bottom edge, that one a side edge, these three
corners are to meet at a shared vertex. Nothing is specified about how
the material between those regions is to be folded, or in what order, or
with what crease pattern. It is left free to buckle, pleat, invert,
overlap, or crumple, provided that whatever it does, the marked
relations end up satisfied.

The Crinkle-Cut Supercube is the constructive example around which this
monograph is built. Three non-contiguous squares were marked on a single
large sheet, at some distance from one another and with no shared edge.
Each square carried edge annotations indicating which of its sides
should occupy the bottom, and which the side, of a finished cube. The
sheet was then folded, without a predetermined crease pattern, until
those three marked squares had been brought together as three mutually
adjacent faces sharing a common corner. The remainder of the sheet,
considerably larger in area than the three faces it was made to support,
was not cut away. It was pushed, pleated, and crumpled into the interior
of the resulting box, where it remains, hidden but present, as the
material cost of the exterior that is visible.

This is not a polished folding technique, and it is not offered as one.
The finished object is visibly irregular. Its value lies elsewhere: it
demonstrates that a target exterior can be specified more sparsely than
a complete crease pattern requires, and that the gap between a sparse
specification and a complete physical object can be closed by permitting
unconstrained material to absorb whatever geometric surplus the
specification does not resolve. What is prescribed is preserved. What is
not prescribed is not thereby prevented; it is simply free, and the
sheet finds \emph{some} way to accommodate it.

The central claim developed across the monograph is the following:

\begin{quote}
Constraint-directed infolding preserves selected surface distinctions by
transferring unresolved geometric surplus into hidden interior
structure.
\end{quote}

Everything that follows is an attempt to make this claim precise, to
give it a formal vocabulary, to trace its consequences for what a
finished boundary can and cannot disclose about its own history, and to
ask how far a strategy discovered by hand with a single sheet of paper
generalizes to other systems in which a legible exterior is achieved
atop a much larger, much less legible interior.

\hypertarget{scope-and-disclaimers}{%
\section{Scope and disclaimers}\label{scope-and-disclaimers}}

Three things this monograph does not attempt should be stated plainly
before the argument begins.

It does not offer a universal folding theorem. The claims advanced here
are existence claims for a demonstrated class of constructions, not
general theorems about which boundary specifications are realizable from
which sheets. Where a stronger conjecture is stated, it is marked as a
conjecture and not treated as established.

It does not offer a complete physical theory of crumpling. Crumpling is
treated throughout as a name for a family of admissible deformations
whose exact microstructure is left underdetermined by design, not as a
phenomenon whose mechanics this monograph resolves. Existing work in the
physics of crumpled sheets is drawn on where relevant but is not the
subject of the monograph.

It does not offer a literal model of cortical development. A structural
analogy is developed in Part V between the paper construction and the
folding of the cerebral cortex, because both involve an area-rich
surface accommodating itself to a volume-limited environment. That
analogy is developed carefully and its limits are stated explicitly, in
the same chapter in which it is introduced, precisely because the
temptation to overstate it is real and the biological literature does
not support a literal reading.

\hypertarget{how-the-argument-proceeds}{%
\section{How the argument proceeds}\label{how-the-argument-proceeds}}

The monograph has three movements.

The first movement, comprising Part I, is constructive. It reconstructs
the original experiment in enough detail to establish exactly what was
demonstrated, and it identifies precisely why existing origami notation,
built around operations and crease patterns, cannot describe what the
experiment did. This motivates the introduction of a fourth descriptive
level, above operations, creases, and shapes: the level of constraints.

The second movement, comprising Parts II and III, is formal. A
vocabulary of marked regions, local frames, boundary conditions, and a
search-like infolding operator is developed in enough detail to state
the Supercube's specification precisely. This formal apparatus is then
used to show why a satisfied specification cannot, by inspection of the
exterior alone, disclose the interior history that satisfied it. This is
the monograph's central theoretical result, and it follows directly from
the formalism rather than being asserted independently of it.

The third movement, comprising Parts IV through VI, extends the argument
outward: to the embodied, externally-scaffolded character of the search
process itself; to a carefully bounded structural analogy with cortical
folding; and finally to a general principle, stated with appropriate
hedges, about when spatially separated marked regions can be brought
into a shared exterior configuration without cutting the sheet that
connects them.

The finished Supercube is, in the end, treated as a witness rather than
as a specimen. It does not prove that every boundary specification is
realizable. It proves that at least one nontrivial specification was,
and it does so in a form that can be inspected, unfolded, and reasoned
about directly. The rest of the monograph is an attempt to say, as
precisely as possible, what that single physical fact actually shows.

\part{The Constructive Problem}

\hypertarget{chapter-1.-a-cube-without-a-net}{%
\chapter{A Cube Without a
Net}\label{chapter-1.-a-cube-without-a-net}}

\hypertarget{the-ordinary-cube-net-as-a-solved-problem}{%
\section{The ordinary cube net as a solved
problem}\label{the-ordinary-cube-net-as-a-solved-problem}}

Before describing what the Crinkle-Cut Supercube did, it is worth being
precise about what an ordinary folded cube does not have to do.

A standard cube net is six squares arranged in the plane so that each
pair of squares meant to become adjacent faces already shares an edge in
the unfolded layout. The classic cross-shaped net is one such
arrangement among eleven possible ones, but every valid net shares the
same property: the planar adjacency of its squares has already been
chosen to match the adjacency the finished cube requires. Folding along
the marked edges does not discover this adjacency. It executes it. The
hard part of the problem, choosing which squares should touch which, has
already been solved on paper before a single fold is made.

This is worth stating because it is easy to mistake the appearance of a
folding task for its actual content. A net looks like raw material
waiting to be shaped, but it is closer to a solution waiting to be
unpacked. The folder's contribution is real but narrow: fold accurately
along the given lines, in a sensible order, and the geometry follows
automatically because it was arranged to follow automatically.

\hypertarget{a-different-starting-point}{%
\section{A different starting
point}\label{a-different-starting-point}}

The Crinkle-Cut Supercube began from a different kind of specification,
one that did not solve the adjacency problem in advance.

A single large sheet, substantially bigger than the six-square area a
minimal net would require, was the starting material. Three squares were
marked on this sheet, and critically, the three squares were not
adjacent to one another and did not share edges. They sat at separated
locations across the sheet, connected only through the intervening
material that lay between them.

Each square was annotated along two of its edges to indicate the role
that edge was meant to play in the finished object. One edge was marked
to become a bottom edge of the eventual cube; another was marked to
become a side edge. A third mark, shared conceptually across the three
squares even though it appeared as a separate small annotation on each,
indicated that a particular corner of each square was meant to become
the same corner of the finished object; the three squares were to meet,
eventually, at one shared exterior vertex.

Nothing else about the sheet was specified. No crease pattern connected
the three squares to one another. No instructions indicated how the
intervening material should be treated, in what order the folds should
proceed, or what shape that material would end up taking. Only the
exterior relations among the three marked squares were declared: which
edges belonged where, and which corners were to coincide.

\hypertarget{what-the-marks-meant}{%
\section{What the marks meant}\label{what-the-marks-meant}}

It is worth dwelling on what kind of information the blue markings
actually carried, because it is easy to mistake them for crease
instructions and they were not that.

A crease instruction says: fold here, in this direction. The blue marks
on the Supercube's three squares said something different: this edge,
wherever it ends up and however it gets there, must be the bottom edge;
this other edge must be the side edge; this corner must coincide with
these two other corners. The marks were coordinate annotations, not
motion instructions. They fixed a target relation among features of the
sheet without fixing the path by which that relation would be reached.

This distinction turns out to matter a great deal, and later chapters
will develop it formally. For now it is enough to notice that the three
marked squares functioned like local frames: each one carried its own
small coordinate system, indicating which of its edges corresponded to
which role in the eventual object, independent of where that square
happened to sit on the original sheet or how far it was from the other
two.

\hypertarget{the-construction}{%
\section{The construction}\label{the-construction}}

The construction proceeded, as closely as the surviving photographs and
the constructor's account allow it to be reconstructed, through four
rough stages.

\textbf{Marking.} The three squares were identified on the flat sheet
and their edges annotated as described above. At this stage the sheet
was still entirely flat, and the three squares were visibly separated by
substantial intervening area, some of it many times the area of the
squares themselves.

\textbf{Approximate alignment.} The sheet was folded, without a
predetermined sequence, in whatever way brought the three marked squares
closer to one another and closer to the relative orientation their
annotations demanded. This stage involved a great deal of trial: a fold
that brought two squares closer together might leave them misoriented
relative to each other, requiring a different fold, or an additional
one, to correct the orientation without losing the proximity already
gained.

\textbf{Corner formation.} At some point in this process, the three
squares arrived at a configuration in which their marked corners could
be brought together at a single shared vertex, with their bottom and
side edges correctly oriented relative to one another. This was the
decisive moment in the construction. Once achieved, it gave the object a
partial cube frame: three faces, correctly adjacent, correctly oriented,
sharing a corner, even though the object as a whole was still far from
resembling a finished cube.

\textbf{Infolding of the remainder.} The material that had not been
selected as one of the three faces, still connected to all three of them
and still comprising the great majority of the sheet's total area, was
folded, pleated, and where necessary crumpled inward, so that it came to
occupy the interior of the box the three faces now defined. This
material was not trimmed or discarded at any point. It was compressed
and redirected until it fit inside the boundary that the three marked
faces had established, becoming the hidden bulk of the finished object.

The photographs of this sequence show exactly this progression: a flat
sheet with a crease pattern and blue markings visible on regions
destined to become exterior faces; a partially folded intermediate state
in which those regions are beginning to approach one another; a
recognizable, if visibly imperfect, cube with the three marked squares
legible as its faces; and, in a separate but related sequence, the same
sheet carried further, through a series of increasingly dense crumples,
into shapes that are no longer cube-like at all, while the same marked
regions remain identifiable within them. This second sequence is
discussed at length in later chapters, since it demonstrates that the
same marked sheet admits more than one admissible resolution and that
the cube was one outcome among a family of outcomes compatible with the
same declared constraints.

\hypertarget{what-was-and-was-not-discarded}{%
\section{What was, and was not,
discarded}\label{what-was-and-was-not-discarded}}

The distinction between the three selected faces and everything else
deserves to be stated as sharply as possible, because the entire
argument of the monograph depends on it.

The unselected material was not waste. It was not treated as scrap to be
trimmed away once its purpose, holding the three squares in their
original flat positions, had been served. It remained physically present
throughout, attached to all three faces at every stage, and it remains
present inside the finished object. What changed was not its existence
but its visibility: it went from being spread out and legible on the
flat sheet to being folded, layered, and hidden inside a boundary that
no longer displays it.

This is the first appearance of a distinction that recurs throughout the
monograph, between a substrate, a realized exterior, and a hidden
remainder. The substrate is the whole sheet. The realized exterior is
the small fraction of it, in this case three squares, that was selected
in advance and that remains legible in the finished object. The hidden
remainder is everything else: present, load-bearing in the sense that it
gives the object its bulk and its stability, but not designed to be
looked at.

\hypertarget{what-the-experiment-demonstrates-precisely}{%
\section{What the experiment demonstrates,
precisely}\label{what-the-experiment-demonstrates-precisely}}

It is important to state the result of this experiment narrowly, because
it is tempting to overstate it.

The Crinkle-Cut Supercube does not demonstrate that any assignment of
three non-contiguous squares to face roles on any sheet can be folded
into a coherent cube corner. It demonstrates one instance. What can be
said with confidence is this:

\begin{quote}
At least one non-contiguous, locally oriented face assignment can be
realized as a coherent box boundary when the intervening surface is
permitted to infold.
\end{quote}

This is a constructive existence claim, and its force comes entirely
from the fact that it is constructive. No argument by symmetry, no
appeal to general folding theorems, and no computational search
established that this configuration was achievable. A physical sheet,
folded by hand, established it, by actually reaching the required
configuration. The object itself is the proof, in the sense that its
existence settles the only question that was actually at issue: whether
the declared relations among the three marked squares were realizable at
all, given a sheet large enough and permission to infold everything
else.

What the experiment leaves open, and what the rest of the monograph is
largely occupied with, is how far this single instance generalizes. Does
the same strategy work for four squares instead of three? For different
relative placements, different orientations, different ratios of
selected area to total sheet area? Does it work reliably, or did this
particular configuration happen to be an easy one? These are the
questions that motivate the formal apparatus developed in Part II, and
they cannot be answered by staring harder at the single artifact
described in this chapter. They require a vocabulary in which the
artifact's construction can be described precisely enough to be varied,
tested, and reasoned about beyond the one sheet that was actually
folded.

Chapter 2 develops the reason such a vocabulary does not already exist
within conventional origami description, and identifies exactly where
the descriptive gap lies.

\hypertarget{chapter-2.-why-existing-origami-notation-is-insufficient}{%
\chapter{Why Existing Origami Notation Is
Insufficient}\label{chapter-2.-why-existing-origami-notation-is-insufficient}}

\hypertarget{what-a-fold-diagram-is-for}{%
\section{What a fold diagram is
for}\label{what-a-fold-diagram-is-for}}

Origami diagramming, in the form most widely used today, descends from
the notation system introduced by Akira Yoshizawa and later standardized
and extended by Samuel Randlett and Robert Harbin. Its vocabulary is by
now familiar even to casual folders: a dashed line for a valley fold, a
dash-dot line for a mountain fold, arrows indicating the direction and
sometimes the destination of a fold, and further symbols for operations
such as reversing a fold, sinking a point inward, rotating the model, or
inflating a pocket.

This notation answers a specific question, asked at a specific moment in
a construction: what should the folder do next? It is addressed to a
reader who has the model in hand, at a known intermediate state, and who
needs to know which edge to bring to which, in which direction, before
turning the page. The notation succeeds by being unambiguous at exactly
this local level. It does not need to explain why the fold is being
made, what larger structure it contributes to, or what would happen if a
different fold were made instead. It is an instruction, not an
explanation, and it is a good one precisely because instructions of this
kind are what a folder executing a known design actually needs.

\hypertarget{what-a-crease-pattern-is-for}{%
\section{What a crease pattern is
for}\label{what-a-crease-pattern-is-for}}

A crease pattern generalizes this a step further. Rather than a sequence
of individual instructions, it presents the complete set of fold lines,
mountain and valley, superimposed on the flattened sheet, from which the
entire model can in principle be folded at once, provided the folder can
determine a workable order for collapsing the pattern.

Crease patterns support a substantial body of mathematics: theorems
about which patterns can be folded flat at all, results connecting local
vertex configurations to global flat-foldability, techniques for
designing a crease pattern to produce a specified silhouette, and
methods for verifying that a proposed folding sequence does not require
the paper to pass through itself. This mathematics is powerful, and it
answers a different question than a step-by-step diagram does: not what
should the folder do next, but what is the complete geometric
relationship between this flat sheet and this folded form?

Both notations, the sequential diagram and the global crease pattern,
share an assumption that is worth making explicit because the
Crinkle-Cut Supercube violates it. Both assume that the relevant crease
structure, whether presented incrementally or all at once, is knowable
and statable in advance, and that describing the fold is largely
equivalent to describing the object. Once the creases are given, the
object is determined, up to the ordinary variability of hand execution.

\hypertarget{practices-that-already-relax-this-assumption}{%
\section{Practices that already relax this
assumption}\label{practices-that-already-relax-this-assumption}}

This assumption has already been relaxed in several existing origami
practices, and it is worth being clear about how far those practices go,
so that what remains undescribed after them can be located precisely.

Wet folding introduces controlled curvature into what would otherwise be
flat facets, using dampened paper's plasticity to hold gentle curves
rather than sharp creases. The resulting surface is not fully specified
by a crease pattern in the traditional sense, since the exact curvature
achieved depends on how the paper is shaped by hand while damp, but the
target form is still generally known in advance and the folder is still
working toward a predetermined shape.

Curved folding extends this further, using continuous curved creases
whose interaction produces developable surfaces that are difficult to
specify by discrete fold lines alone. This mathematics is more demanding
than flat-crease theory, but it remains a theory of how a known target
surface arises from a known family of creases.

Crumpling, treated as a technique in its own right rather than as an
accident to be avoided, appears in sculptural and textural origami,
where an artist may crumple paper deliberately to produce organic,
non-planar surface texture as part of a larger design. Here the exact
microstructure of the crumple is genuinely left unspecified. But the
crumpled region is typically decorative or textural. It is not usually
required to satisfy a precise exterior relation to some other,
separately specified part of the model. Its underdetermination is a
feature of its appearance, not a device for realizing a boundary
constraint that would otherwise be unreachable.

Each of these practices weakens the connection between crease
specification and final form in a particular, limited way. None of them
provides a vocabulary for the situation the Supercube presents: several
widely separated regions of a sheet, each carrying its own oriented
target relation to the finished exterior, with the entire remainder of
the sheet left free to do whatever is necessary, structurally rather
than decoratively, to bring those regions into their required relation.

\hypertarget{the-question-none-of-these-notations-ask}{%
\section{The question none of these notations
ask}\label{the-question-none-of-these-notations-ask}}

It is possible to state the gap directly. Sequential diagrams answer:
what should the folder do next? Crease patterns answer: along which
lines must the sheet fold to produce this form? Wet folding and curved
folding answer variants of the same question, extended to continuous
curvature. Textural crumpling answers: how can this region's appearance
be varied without a fixed target? None of them answers the question the
Supercube's construction was actually organized around:

\begin{quote}
Which relations must hold at the final boundary, regardless of the exact
crease network through which they are realized?
\end{quote}

This is a question about the object's exterior, posed independently of
any particular route to achieving it. It does not ask which creases to
make. It asks what must be true when the folding is finished, leaving
the folding itself, and everything about it that does not bear on that
final truth, unconstrained. A crease pattern can be seen as one
particular, fully-specified answer to a question of this kind, in the
special case where every degree of freedom has been resolved in advance.
The Supercube's specification is a much sparser answer to the same
general kind of question, one that resolves only the exterior relations
among three regions and leaves everything else open.

\hypertarget{four-descriptive-levels}{%
\section{Four descriptive levels}\label{four-descriptive-levels}}

It is useful to organize this discussion around four distinct levels at
which a paper construction can be described, moving from the most local
and procedural to the most global and declarative.

At the \textbf{operation level}, description consists of individual
actions: fold this edge to that one, reverse this flap, sink this point.
This is the level of the sequential diagram. It is addressed to the
hands, one instruction at a time.

At the \textbf{crease level}, description consists of the complete
pattern of fold lines on the flattened sheet, without necessarily
specifying an order of execution. This is the level of crease-pattern
mathematics. It is addressed to the sheet as a whole, but still in terms
of where it must fold.

At the \textbf{shape level}, description consists of the target
three-dimensional form itself: a cube, a crane, a tessellated surface.
This is the level at which a finished model is usually named and
recognized. It says what the object is, without saying how it got that
way.

At the \textbf{constraint level}, description consists of a set of
relations that selected features of the finished object must satisfy:
this edge here, this corner there, these three regions mutually
adjacent, this much material relegated to the interior. This level does
not specify operations, does not specify a complete crease pattern, and
does not even specify the complete shape. It specifies only the subset
of the shape's structure that has been chosen to matter, and it leaves
everything else, including the exact crease network by which those
chosen relations come to be satisfied, as free variables to be resolved
during construction rather than before it.

Existing origami description operates almost entirely at the first two
levels, with the third level, the target shape, usually serving as an
informal label attached to a complete crease-level specification rather
than as an independent description in its own right. The Crinkle-Cut
Supercube's specification belongs to the fourth level. Three squares and
their edge and corner annotations constitute a constraint-level
description: sparse, declarative, and silent about the crease network
required to satisfy it. The remaining chapters of Part II are devoted to
giving this fourth level a formal vocabulary of its own, since none of
the existing levels can be stretched to cover it without losing the very
property, the deliberate underdetermination of everything not explicitly
marked, that makes the Supercube's construction what it is.

\hypertarget{reproducibility-reconsidered}{%
\section{Reproducibility
reconsidered}\label{reproducibility-reconsidered}}

One further consequence of working at the constraint level is worth
drawing out before moving on, because it changes what it would mean for
someone else to reproduce the Crinkle-Cut Supercube.

A conventional diagram aims at a strong form of reproducibility: a
second folder, following the same instructions, should arrive at very
nearly the same crease network and the same final layering as the
original. Fidelity to the instructions is fidelity to the object.

A constraint-level specification supports, and arguably only supports, a
weaker but more interesting form of reproducibility. A second folder
given the same three marked squares and the same edge and corner
annotations, but no information about how the original sheet was
actually folded, might arrive at a finished cube corner by an entirely
different sequence of folds, with an entirely different interior crease
network and layering, while still satisfying the same declared
relations. Both constructions would be, in a sense to be made precise in
Chapter 6, equivalent: not because they share a folding history, but
because they share a boundary.

This is not a weaker notion of success. It is a different one, and it is
the notion the rest of this monograph will be built around. Two
realizations count as satisfying the same specification when they agree
on what was declared to matter, whatever else may differ between them.
Chapter 3 begins the work of stating what, exactly, is declared to
matter, and how.

\part{A Formal Language of Infolding}

\hypertarget{chapter-3.-marked-surfaces-and-local-frames}{%
\chapter{Marked Surfaces and Local
Frames}\label{chapter-3.-marked-surfaces-and-local-frames}}

\hypertarget{the-sheet-as-a-material-surface}{%
\section{The sheet as a material
surface}\label{the-sheet-as-a-material-surface}}

The formal apparatus developed in this and the following two chapters is
built to describe exactly the situation Chapter 1 reconstructed: a
single connected sheet, most of it destined to disappear from view, a
small number of regions on it singled out in advance, and a target
relation among those regions that the finished object must satisfy.

Let \(P\) denote the sheet. At the coarsest level of description, \(P\)
can be treated as a connected, bounded two-dimensional surface,
initially lying flat in the plane. This idealization is useful for
stating relations precisely, but it should not be mistaken for a claim
that the sheet's material properties are irrelevant. Real paper has
thickness, stiffness, grain direction, and a limited capacity to bend
before it creases plastically rather than elastically; it resists
self-intersection and cannot pass through itself. These properties
govern what is physically achievable and are reintroduced wherever the
argument depends on them, particularly in Chapter 6's treatment of
surplus area and Chapter 9's treatment of obstructions. For the purpose
of stating what is being specified, however, it is enough to treat \(P\)
as a surface whose points can be tracked through a deformation.

\hypertarget{selected-regions-and-the-infoldable-remainder}{%
\section{Selected regions and the infoldable
remainder}\label{selected-regions-and-the-infoldable-remainder}}

A finite collection of regions is selected on \(P\) in advance:

\[E = \{E_1, E_2, \ldots, E_n\}\]

These are the regions intended to remain legible on the finished
object's exterior. In the Supercube, \(n = 3\), and each \(E_i\) is one
of the three marked squares.

The complement of this selection,

\[I = P \setminus E\]

is the infoldable remainder: every part of the sheet not singled out for
exterior legibility.

The word ``infoldable'' is chosen deliberately over alternatives like
``excess'' or ``waste,'' because those alternatives suggest that \(I\)
is material to be gotten rid of, and it is not. Nothing about \(I\) is
discarded, cut, or treated as disposable. What distinguishes \(I\) from
\(E\) is only that its final geometry has not been declared in advance.
It is free where \(E\) is fixed. This is a statement about how much has
been specified, not a statement about the region's importance to the
finished object; as Chapter 6 will argue at length, \(I\) typically
constitutes the great majority of the object's material bulk even though
it constitutes none of its declared exterior.

\hypertarget{local-frames}{%
\section{Local frames}\label{local-frames}}

Each selected region carries more than just its own extent. It carries
an orientation: a declaration of which of its features correspond to
which roles in the target exterior.

Write a framed region as a pair:

\[F_i = (E_i, o_i)\]

where \(o_i\) is the orientation data attached to \(E_i\). For a square
region, this orientation data might distinguish a bottom edge from a
side edge from a top edge, or equivalently specify a distinguished
corner together with the two edge directions incident to it. The precise
content of \(o_i\) depends on the shape of \(E_i\) and on how many of
its features are given a target role; nothing in the formalism requires
\(E_i\) to be a square, or requires every edge of \(E_i\) to be labeled.

In the Supercube, the physical blue marks are best understood as a
direct, material instantiation of this orientation data. They were not
drawn to indicate where the sheet should fold. They were drawn to
indicate which edge of each marked square was to become which edge of
the finished cube, independent of the square's position on the original
flat sheet or its distance from the other two marked squares. A local
frame, in other words, is portable in exactly the sense the marked
squares needed to be: it travels with the region through whatever
deformation eventually brings that region into its target position, and
it is legible on the region wherever that region ends up.

\hypertarget{what-a-local-frame-does-not-specify}{%
\section{What a local frame does not
specify}\label{what-a-local-frame-does-not-specify}}

It is worth being precise about the boundary of what a local frame
declares, because the formalism depends on this boundary being sharp.

A local frame \(F_i = (E_i, o_i)\) specifies the identity of a region
and the roles assigned to some of its features. It does not specify
where \(E_i\) currently sits relative to the other selected regions. It
does not specify how \(E_i\) is to be moved, folded, or reoriented in
order to bring \(o_i\) into agreement with the corresponding roles on
the target object. And critically, it says nothing at all about the
material outside \(E_i\), including the material immediately adjacent to
it. A local frame is a statement about one region, made independently of
every other region and independently of the deformation that will
eventually be required to realize it.

This independence is what makes it possible to declare relations among
widely separated regions without having to say anything about the
intervening sheet. The three squares of the Supercube could be marked
without reference to one another and without reference to the material
lying between them, because each mark was a self-contained statement
about one region's intended role. The relations \emph{among} the three
frames, which edges must meet which, which corners must coincide, are a
separate layer of specification, introduced in Chapter 4. A local frame,
on its own, only ever describes a single region's target orientation.

\hypertarget{edge-roles-and-partial-labeling}{%
\section{Edge roles and partial
labeling}\label{edge-roles-and-partial-labeling}}

Within a local frame, individual edges (or, more generally, individual
boundary features) can be assigned symbolic roles. A short,
non-exhaustive vocabulary suffices for most purposes:

\begin{itemize}
\tightlist
\item
  \(B\), for an edge intended to become a bottom edge of the target
  object;
\item
  \(S\), for an edge intended to become a side edge;
\item
  \(T\), for an edge intended to become a top edge;
\item
  \(J_k\), for an edge intended to join a specific other labeled edge,
  indexed by \(k\);
\item
  \(V_k\), for an edge incident to a designated target vertex, indexed
  by \(k\);
\item
  \(X\), for an edge explicitly prohibited from exterior exposure,
  should such a prohibition be needed.
\end{itemize}

A framed region can then be written compactly. For instance, one of the
Supercube's three squares might be written:

\[F_1 = [E_1;\ B = e_1,\ S = e_2]\]

indicating that edge \(e_1\) of region \(E_1\) is assigned the bottom
role and edge \(e_2\) is assigned the side role, with the remaining
edges of \(E_1\) left unlabeled.

That last clause, the remaining edges left unlabeled, is not an
oversight to be corrected in a more careful specification. It is the
specification. A local frame need not assign a role to every edge of its
region. Leaving an edge unlabeled is a positive declaration that the
final orientation or adjacency of that edge is not being constrained,
not a gap that a more thorough author would have filled in. This is what
is meant, throughout this monograph, by a \textbf{partial
specification}: not an incomplete draft of a complete one, but a
deliberately incomplete final statement of exactly those relations the
constructor has chosen to care about.

\hypertarget{why-partiality-is-load-bearing}{%
\section{Why partiality is
load-bearing}\label{why-partiality-is-load-bearing}}

It is worth pausing on why this partiality matters, rather than treating
it as an incidental economy of notation.

A specification that labeled every edge of every region, and
additionally specified the exact relation of every point of \(I\) to
every other point, would no longer be a constraint-level description in
the sense of Chapter 2. It would have collapsed back into something
closer to a complete crease pattern: a full account of the finished
object's geometry, articulated region by region rather than crease by
crease, but no less total in its coverage. Total specification of this
kind is not wrong, but it is a different project. It reintroduces
exactly the assumption Chapter 2 identified as the limiting feature of
existing notation, that describing the object is equivalent to
describing its complete geometry, and it forecloses the very freedom
that allowed the Supercube's surrounding material to find its own
accommodation.

The formalism's usefulness depends on keeping the labeled portion small
relative to the unlabeled portion. A specification that declares three
edges and one shared vertex, out of a sheet with area many times larger
than the three regions those edges belong to, constrains a
correspondingly enormous space of possible deformations down to whatever
subset of that space happens to satisfy the three declared relations.
The interesting content of the theory, developed across the remaining
chapters of Part II and the whole of Part III, concerns what can be said
about that subset, how large or small it is, what its members have in
common beyond the declared relations themselves, and what physical
process might be relied upon to locate a member of it. None of this
content is visible if the specification is allowed to grow until it
covers the whole sheet.

\hypertarget{sparse-specification-as-the-chapters-central-idea}{%
\section{Sparse specification as the chapter's central
idea}\label{sparse-specification-as-the-chapters-central-idea}}

The formal objects introduced in this chapter, the substrate \(P\), the
selected regions \(E\) and their complement \(I\), and the local frames
\(F_i\) with their partially labeled edge roles, are deliberately
minimal. They say only as much as the Supercube's actual construction
said: which regions matter, and which of their features are assigned
target roles. Everything else about the sheet's eventual geometry is
left open by construction, not by omission.

This is the sense in which the Supercube's specification is
\textbf{sparse}: a small, explicitly bounded set of declared relations
governing a much larger surface, with no claim made, and no claim
needed, about anything the specification does not mention. Chapter 4
extends this vocabulary from single regions to relations among regions,
stating formally what it means for two or more framed regions to be
required to meet, align, or coincide at the finished object's boundary,
and introduces the graphic notation by which such a specification can be
recorded on the sheet itself, in the same spirit as the original blue
marks.

\hypertarget{chapter-4.-boundary-grammar-and-infolding-notation}{%
\chapter{Boundary Grammar and Infolding
Notation}\label{chapter-4.-boundary-grammar-and-infolding-notation}}

\hypertarget{from-single-frames-to-relations-among-frames}{%
\section{From single frames to relations among
frames}\label{from-single-frames-to-relations-among-frames}}

Chapter 3 gave a vocabulary for describing one selected region at a
time: its identity, its extent, and the roles assigned to some of its
edges. That vocabulary is not yet sufficient to state the Supercube's
specification, because the specification's real content lies in the
relations \emph{among} the three marked squares, not in any one square
considered alone. This chapter introduces that relational layer,
together with a compact graphic notation for recording it directly on a
sheet.

Let \(\Omega\) denote the effective volume of the object that results
from a given deformation, and let \(\partial\Omega\) denote its apparent
exterior boundary. A \textbf{realization} is a deformation \(f\)
carrying the marked sheet \(P\) into a folded state that satisfies a
declared set of constraints \(C\). The remainder of this chapter defines
the constraint types that make up \(C\).

\hypertarget{exposure}{%
\section{Exposure}\label{exposure}}

The simplest constraint asserts that a selected region ends up visible
on the finished exterior:

\[f(E_i) \subset \partial\Omega\]

within whatever tolerance the construction's purpose demands. An exact
reading of this condition would require every point of \(E_i\) to lie
precisely on the boundary surface; a more forgiving reading, appropriate
to a hand-folded object like the Supercube, allows \(E_i\) to be
substantially, though not perfectly, exposed, with minor overlap or
partial occlusion at its edges tolerated. Chapter 6 returns to the
question of how such tolerance should be measured; for now it is enough
to note that an exposure condition is a claim about where a region ends
up, not about its exact shape once it gets there.

\hypertarget{adjacency}{%
\section{Adjacency}\label{adjacency}}

An adjacency condition requires a specified edge of one region to meet a
specified edge of another in the finished exterior:

\[A(E_i:e_a,\ E_j:e_b)\]

This notation reads: edge \(e_a\) of region \(E_i\) must be adjacent, in
the realized object, to edge \(e_b\) of region \(E_j\). Adjacency can be
directed, when the orientation of the meeting matters, distinguishing an
aligned meeting from a reversed one, and it can further distinguish a
flush meeting from a merely incident one, in which the two edges touch
without running parallel along their full length. The Supercube's
construction relied on adjacency conditions of exactly this kind to
specify that the bottom edge of one marked square was to meet the side
edge of another.

\hypertarget{vertex-gluing}{%
\section{Vertex-gluing}\label{vertex-gluing}}

Where three or more frames are required to meet not merely edge-to-edge
but at a single common point, a vertex-gluing condition is used:

\[G(F_1, F_2, F_3) = V\]

This states that the designated corners of the three named frames must
become incident to one shared exterior vertex \(V\) in the realized
object. This was the single most consequential condition in the
Supercube's specification. Chapter 1 described the moment at which this
condition first became satisfied as the decisive turn in the
construction: once the three marked corners could be brought together at
one point with their associated edges correctly oriented, the object
possessed a coherent partial cube frame, even though the great majority
of the sheet remained unresolved. A vertex-gluing condition is, in this
sense, disproportionately powerful relative to its brevity: it ties
together conditions on several frames at once, and its satisfaction
tends to lock a large amount of the surrounding geometry into place as a
side effect.

\hypertarget{orientation}{%
\section{Orientation}\label{orientation}}

A region may satisfy every adjacency and vertex condition demanded of it
while still being incorrectly rotated or reflected relative to its
intended role, a bottom edge occupying a side position, for instance,
even where the correct edges happen to touch. An orientation condition
rules this out by requiring the realized deformation to carry a frame's
declared orientation into agreement with its target orientation:

\[D f_i(o_i) \approx O_i\]

The notation here is not meant to invoke a literal derivative or to
require smooth differentiability of \(f\) everywhere; the sheet's actual
deformation includes sharp creases at which no such derivative would be
classically defined. What the notation expresses is the preservation of
the directional roles declared in \(o_i\): that whatever the deformation
does elsewhere, the marked bottom edge ends up oriented as a bottom
edge, the marked side edge as a side edge, and so on, up to the
tolerance the construction allows.

\hypertarget{relegation}{%
\section{Relegation}\label{relegation}}

The infoldable remainder \(I\), defined in Chapter 3, is not left
entirely unconstrained. It is subject to one standing condition, which
can be stated as a relegation constraint:

\[f(I) \subset \Omega \cup H\]

where \(H\) denotes the collection of hidden boundary layers, seams,
overlaps, and interior pockets the realized object may contain. This is
not a claim that the free material occupies a literal, zero-thickness
mathematical interior; a folded sheet has thickness, and its layers
press against one another and against the interior surfaces of the
exposed faces. The condition should be read practically: the free
material must end up tucked inside the apparent envelope of the object,
or hidden beneath its designated exterior faces, rather than protruding
as unaccounted-for exterior surface of its own. Relegation is what
allows the object's declared exterior, in the Supercube's case just
three squares, to remain the only legible surface of an object built
from a sheet many times larger than those three squares combined.

\hypertarget{closure}{%
\section{Closure}\label{closure}}

A final constraint type governs how complete, stable, or functionally
container-like the finished object must be. Closure is deliberately left
as a family of conditions rather than a single formula, since what
counts as adequate closure depends entirely on the construction's
purpose. Possible closure conditions include: that the boundary contains
no large unintended openings; that the relegated material remains
contained without external support once the folder's hands are removed;
that the object holds its shape for some specified duration; or, at the
least demanding end, only that the vertex-gluing and adjacency
conditions above remain satisfied, with no further requirement on the
object's overall rigidity or permanence.

The Supercube's own closure was modest by this reckoning: the three
marked faces held their adjacency and vertex relations, and the bulk of
the sheet was contained within the resulting boundary, but the object
was not rigid, was not sealed against all gaps, and would not have
survived indefinite handling without some of its relations loosening.
This is treated in this monograph not as a shortfall but as an accurate
report of exactly which conditions the construction was actually able to
secure, a report that later chapters, particularly Chapter 9's
discussion of protected and sacrificial distinctions, will show was
itself a meaningful outcome rather than a simple failure to finish the
job.

\hypertarget{the-boundary-grammar-as-a-whole}{%
\section{The boundary grammar as a
whole}\label{the-boundary-grammar-as-a-whole}}

Together, exposure, adjacency, vertex-gluing, orientation, relegation,
and closure constitute what this monograph calls a \textbf{boundary
grammar}: a set of relation types by which a finished object's exterior
can be specified without specifying the deformation that produces it.
The Supercube's full specification can now be written compactly,
gathering the pieces introduced across this chapter and the last:

\begin{gather*}
\text{INFOLD}(P \setminus E \mid \text{Exposure}(E),\ \text{Adjacency}(A),\\
\text{Orientation}(O),\ \text{Vertex}(V),\ \text{Closure}(K))
\end{gather*}

The left-hand argument, \(P \setminus E\), names the free material to be
searched over. The right-hand arguments name the boundary grammar
conditions that free material's eventual configuration must satisfy, or
must not prevent the rest of the object from satisfying. Chapter 5 takes
up the \(\text{INFOLD}\) operator itself, and the search process it
denotes, in detail. The remainder of this chapter turns from the
algebraic notation above to a graphic notation suited to marking these
same conditions directly on a physical sheet, in the spirit of the
original blue annotations.

\hypertarget{a-graphic-vocabulary}{%
\section{A graphic vocabulary}\label{a-graphic-vocabulary}}

A workable notation for constraint-directed infolding should be usable
by hand, without requiring a computational model to interpret it, and it
should visually distinguish protected material from free material at a
glance. The following elements are proposed to meet those requirements;
the exact graphic conventions are not claimed to be uniquely correct,
only sufficient, and Appendix A collects them in a single reference
sheet.

\textbf{Protected faces.} A selected region \(E_i\) is marked with a
heavy outline and given a short identifier, \(F_1\), \(F_2\), \(F_3\),
and so on. The heavy outline distinguishes it visually, at a glance,
from the surrounding free material.

\textbf{Free material.} Regions belonging to \(I\), the infoldable
remainder, are left unmarked or given only a light shading. This absence
of marking is itself informative: it signals that the region's internal
crease structure is outside the specification, not merely that it has
not yet been decided how to draw it.

\textbf{Edge-role ticks.} Rather than adopting mountain- and valley-fold
lines for this purpose, which already carry a specific and different
meaning, edge roles are marked with short ticks or letters placed near
the relevant edge: a single solid tick for a bottom-role edge, a double
tick for a side-role edge, and a matching pair of letters, placed near
two edges required to meet, to indicate an adjacency condition between
them.

\textbf{Vertex labels.} A small circled label, such as \(V_1\), repeated
at a corner of each of several selected regions, indicates that those
corners are required to become incident to one shared exterior vertex.
Repetition of the same label across multiple regions is what carries the
vertex-gluing condition; a single unrepeated label carries no such
requirement.

\textbf{Exposure status.} Where a region's exposure is not simply
assumed, a short status label may be attached: \(\text{OUT}\) for a
region intended to be principally exterior, \(\text{IN}\) for one
intended to be principally interior, \(\text{FREE}\) where either status
is acceptable, and \(\text{SEAM}\) for a region permitted to appear only
as a narrow boundary strip or overlap rather than as a full exterior
face.

\textbf{Phase annotations.} Where a construction benefits from being
described in stages, brief phase labels may be added separately from the
geometric marks themselves: \emph{realize \(V_1\)}, then \emph{align the
bottom-edge cycle}, then \emph{relegate the remaining surplus}. These
annotations describe the order in which constraints are to be secured,
not the exact hand motions by which any one of them is secured; they are
closer to a solver's schedule than to a fold-by-fold diagram.

\hypertarget{compatibility-with-conventional-symbols}{%
\section{Compatibility with conventional
symbols}\label{compatibility-with-conventional-symbols}}

None of the notation introduced in this chapter is intended to replace
the existing vocabulary of mountain and valley folds, reversals, sinks,
and the rest. A hybrid diagram is not only permitted but often
desirable: conventional fold symbols can specify a known local
operation, perhaps a reversal used to bring two marked edges into
contact, while the constraint-level symbols introduced here specify what
that operation is being performed in order to preserve or achieve. The
two vocabularies operate at different descriptive levels, in the sense
of Chapter 2, and a complete practical diagram for a constraint-directed
construction may well need both: the constraint level to state what must
be true when the folder is finished, and the operation level to record
whatever specific moves the folder actually found useful along the way.

Chapter 5 turns to the process this notation describes rather than
merely records: the search, denoted \(\text{INFOLD}\), by which a folder
actually locates a deformation satisfying a declared boundary grammar,
and the auxiliary operations, locking, relaxing, and repairing a
constraint mid-construction, that make such a search tractable by hand.

\hypertarget{chapter-5.-the-infold-operator-and-constraint-preserving-search}{%
\chapter{The Infold Operator and Constraint-Preserving
Search}\label{chapter-5.-the-infold-operator-and-constraint-preserving-search}}

\hypertarget{a-different-kind-of-instruction}{%
\section{A different kind of
instruction}\label{a-different-kind-of-instruction}}

The notation developed in Chapter 4 describes what a finished
realization must satisfy. It does not, by itself, describe the process
of arriving at one. This chapter introduces that process, in the form of
a single operator that stands in deliberate contrast to an ordinary fold
instruction.

An ordinary fold instruction, of the kind Chapter 2 associated with the
operation level of description, names a specific action: bring this edge
to that one, in this direction. Executed correctly, it produces one
determinate result. The operator introduced here, written
\(\text{INFOLD}\), does not name a specific action and does not produce
one determinate result. It names a family of permitted actions and a set
of conditions those actions must eventually satisfy, and it leaves the
choice of any particular sequence of actions unresolved.

\[\text{INFOLD}(I \mid C)\]

reads: deform the free region \(I\) through any combination of the
permitted operations, in any order and to any extent, such that the
constraint set \(C\) comes to be satisfied, or is not thereby made
unsatisfiable. \(\text{INFOLD}\) is better understood as a search
instruction than as a fold instruction. It tells the constructor what
must be true at the end without telling them how to get there, in the
same sense that a specification tells an engineer what a bridge must
bear without telling them which member to weld first.

\hypertarget{the-supercubes-instruction-written-out}{%
\section{The Supercube's instruction, written
out}\label{the-supercubes-instruction-written-out}}

Gathering the constraint types introduced in Chapter 4, the Supercube's
full specification can be written:

\begin{gather*}
\text{INFOLD}(P \setminus E \mid \text{Exposure}(E),\ \text{Adjacency}(A),\\
\text{Orientation}(O),\ \text{Vertex}(V),\ \text{Closure}(K))
\end{gather*}

Read in full: deform the material outside the three marked squares, by
any combination of folding, pleating, and crumpling, such that the three
squares remain exposed, their marked edges meet as declared, their
marked orientations are preserved, their marked corners coincide at a
shared vertex, and the resulting object achieves at least the modest
closure the construction actually required. Nothing in this instruction
specifies which fold to make first, how many creases the free material
should acquire, or what shape that material takes once it is no longer
visible. Those questions are left to whatever process actually carries
out the search, which is the subject of the remainder of this chapter
and, from a different angle, of Chapter 7.

\hypertarget{the-permitted-operation-family}{%
\section{The permitted operation
family}\label{the-permitted-operation-family}}

\(\text{INFOLD}\) ranges over a specific, if broad, family of physical
operations. Bending, buckling, pleating, corrugating, sinking, local
inversion, crumpling, nesting, overlapping, wrapping, and the insertion
of one layer beneath another are all permitted. Cutting and adhesive
bonding are not permitted by default; a construction that requires
either should say so explicitly, since both operations change the
sheet's connectivity in ways the rest of this monograph's argument,
particularly Chapter 6's treatment of intrinsic connectivity, depends on
being ruled out unless stated otherwise. Tearing occupies an
intermediate position: depending on the construction's purpose, an
accidental tear may either invalidate the realization outright or be
treated as a separate, weaker category of success, a distinction Chapter
9 returns to when it discusses obstructions to realizability.

\hypertarget{protected-and-sacrificial-geometry}{%
\section{Protected and sacrificial
geometry}\label{protected-and-sacrificial-geometry}}

The asymmetry between the three marked squares and the rest of the
sheet, already familiar from Chapters 3 and 4, can now be stated as an
asymmetry in how much freedom \(\text{INFOLD}\) grants to different
parts of \(I\).

Some geometry is \textbf{protected}: it corresponds directly to a
declared constraint, and \(\text{INFOLD}\) must not sacrifice it in the
course of resolving anything else. The three marked squares' exposure,
their edge roles, and their shared vertex are protected in this sense
throughout the Supercube's construction.

Other geometry is \textbf{sacrificial}: its precise final shape and its
crease history are explicitly not protected, and \(\text{INFOLD}\) is
free to let this geometry take whatever form is convenient for
satisfying the protected constraints. It is important to be precise
about what ``sacrificial'' means here, since the word can misleadingly
suggest destruction. No material is destroyed. What is sacrificed is
specificity: the constructor gives up any claim to control the exact
shape this geometry ends up taking, in exchange for the freedom that
surrender buys elsewhere in the construction. The bulk of \(I\) in the
Supercube is sacrificial in exactly this sense. Its final crumpled form
was never specified, was not predictable in advance, and does not need
to be reproduced by a second constructor working from the same
specification; only the protected relations need to be reproduced.

This asymmetry, high constraint on a small selected exterior, low
constraint on a large hidden connective mass, is the operational heart
of constraint-directed infolding. It is what allows a sparse
specification, in the sense of Chapter 3, to be realized at all: the
freedom given up on the sacrificial majority of the sheet is exactly the
freedom that makes it possible to satisfy demanding relations on the
small protected minority.

\hypertarget{locking-a-relation-once-achieved}{%
\section{Locking a relation once
achieved}\label{locking-a-relation-once-achieved}}

A search process that has to keep every constraint simultaneously at
risk throughout construction is harder to carry out, by hand or
otherwise, than one that can secure a relation and then set it aside.
The Supercube's own construction, as reconstructed in Chapter 1, relied
on securing the vertex-gluing condition first and only then addressing
the remaining material, and this reliance can be given a name.

\[\text{LOCK}(C_k)\]

denotes the decision to treat a currently satisfied constraint \(C_k\)
as provisionally fixed, so that subsequent operations are chosen, or
avoided, specifically to preserve it. For the Supercube, once the three
marked corners had been brought together at a shared vertex, that
relation was locked:

\[\text{LOCK}\big(G(F_1,F_2,F_3) = V\big)\]

and the remaining infolding was carried out under this lock, meaning
that no subsequent manipulation of the free material was permitted to
disturb the vertex once it had been achieved. Combined with the operator
introduced above, the full sequence can be written:

\[\text{INFOLD}\big(I \mid \text{LOCK}(V),\ \text{Closure}(K)\big)\]

This expresses precisely the strategy Chapter 1 described in narrative
form: establish the three-face corner first, treat it thereafter as
fixed, and resolve everything else in whatever way that fixed corner
permits.

\hypertarget{relaxing-a-lock-that-has-become-an-obstruction}{%
\section{Relaxing a lock that has become an
obstruction}\label{relaxing-a-lock-that-has-become-an-obstruction}}

A locked relation, useful as it is for organizing a search, can also
become an obstruction. Material rearranged elsewhere in the sheet may,
at some point, be unable to proceed without disturbing a relation that
was locked earlier. Rather than treating this as a dead end, the
formalism permits a controlled relaxation:

\[\text{RELAX}(C_k, \varepsilon)\]

which permits the previously locked relation \(C_k\) to drift,
temporarily, within a stated tolerance \(\varepsilon\), so that the
surrounding obstruction can be cleared. This is not an abandonment of
the constraint. It is an explicit, bounded loosening, made in service of
eventually satisfying the construction as a whole, and it is paired with
a corresponding operation for restoring the relation once the
obstruction has passed:

\[\text{REPAIR}(C_k)\]

which reasserts \(C_k\), ideally within its original tolerance, once the
surrounding material has been rearranged sufficiently to permit this.

\hypertarget{why-relaxation-and-repair-matter-beyond-bookkeeping}{%
\section{Why relaxation and repair matter beyond
bookkeeping}\label{why-relaxation-and-repair-matter-beyond-bookkeeping}}

It would be possible to regard \(\text{LOCK}\), \(\text{RELAX}\), and
\(\text{REPAIR}\) as mere bookkeeping devices, convenient shorthand for
describing a process that could equally well be described without them.
This understates their significance.

What these three operations jointly capture is that constraint-directed
infolding, as actually practiced by hand, is not a monotonic process in
which each constraint, once satisfied, remains satisfied until the end.
It is a process with local setbacks that are recoverable precisely
because the constraint was never irreversibly fused with the surrounding
sheet; a locked relation, expressed this way, is always understood to be
provisionally locked, subject to exactly the kind of temporary loosening
\(\text{RELAX}\) names. A notation that only recorded successes, and had
no vocabulary for a temporary, bounded, intentional step backward, would
misrepresent how the construction actually proceeded, and would offer no
way to distinguish a principled, recoverable relaxation from an
accidental, possibly unrecoverable failure of the same relation. Chapter
7 develops the embodied side of this same point, treating a blocked fold
not as an error external to the search but as information about the
current boundary of what is reachable; \(\text{RELAX}\) and
\(\text{REPAIR}\) are this chapter's formal counterpart to that later,
more experiential discussion.

\hypertarget{what-chapter-5-has-and-has-not-established}{%
\section{What Chapter 5 has and has not
established}\label{what-chapter-5-has-and-has-not-established}}

This chapter has given the Supercube's construction a precise
operational vocabulary: an operator, \(\text{INFOLD}\), denoting
constrained search over a permitted family of deformations; a
distinction between protected and sacrificial geometry that explains why
such a search is tractable at all; and a small set of auxiliary
operations, \(\text{LOCK}\), \(\text{RELAX}\), and \(\text{REPAIR}\),
that describe how a constraint can be provisionally secured, temporarily
loosened, and subsequently restored over the course of a single
construction.

What it has not established, and what the remaining chapters of the
monograph take up in turn, is any general account of when a given
\(\text{INFOLD}\) instruction is satisfiable at all, how large the space
of satisfying deformations tends to be, or what, if anything, the
members of that space have in common beyond satisfying the same declared
constraints. These are the questions Part III addresses, beginning with
Chapter 6's treatment of the relationship between a sheet's original
connectivity and the exterior adjacency its infolded remainder makes
possible.

\part{Hidden Geometry and Exterior Underdetermination}

\hypertarget{chapter-6.-adjacency-surplus-surface-and-historical-interior}{%
\chapter{Adjacency, Surplus Surface, and Historical
Interior}\label{chapter-6.-adjacency-surplus-surface-and-historical-interior}}

\hypertarget{two-graphs-on-the-same-sheet}{%
\section{Two graphs on the same
sheet}\label{two-graphs-on-the-same-sheet}}

The Supercube's construction can be described by two different graphs,
and the entire theoretical argument of this chapter follows from
noticing that these two graphs need not agree.

The first is the \textbf{sheet graph}, written \(G_0\), whose nodes
correspond to the marked regions and the intervening portions of the
flat, unfolded sheet, and whose edges record which of these regions are
directly connected across the original planar surface. In the Supercube,
\(G_0\) records the fact that the three marked squares are widely
separated, joined only through a large connected mass of unmarked
material, with no two of the three squares sharing an edge anywhere on
the flat sheet.

The second is the \textbf{boundary graph}, written \(G_\partial\), whose
nodes correspond to the regions actually visible on the finished
exterior, and whose edges record which of these regions are adjacent to
one another in the realized object. In the Supercube, \(G_\partial\)
records that the three marked squares, after folding, are mutually
adjacent, meeting pairwise along the edges their annotations declared
and sharing the common vertex their vertex-gluing condition demanded.

The Supercube's construction was organized around producing a specific,
wanted disagreement between these two graphs:

\[G_\partial \neq \text{(the exposed subgraph induced by } G_0\text{)}\]

The three squares are not adjacent in \(G_0\). They are adjacent in
\(G_\partial\). Something has to account for this difference, and it is
worth being precise about what does and does not account for it.

\hypertarget{what-changes-and-what-does-not}{%
\section{What changes, and what does
not}\label{what-changes-and-what-does-not}}

It might seem that producing new adjacency where none existed before
requires cutting the sheet, since cutting is the ordinary way material
connectivity is altered. This is not what happened, and the reason it is
not what happened is the central fact this chapter exists to draw out.

The sheet's intrinsic material connectivity, its status as a single
connected piece of paper with no tears or seams introduced, is unchanged
by the construction. Every point of the original sheet remains connected
to every other point by some continuous path through the material,
exactly as it was before any fold was made. What changes is which of
these connections lie on the exterior boundary of the finished object
and which have been folded away from it.

The two squares that end up adjacent in \(G_\partial\) were never made
adjacent in \(G_0\); no new material connection was created between
them. Instead, the material that used to separate them in the flat
sheet, previously visible and lying between them in the plane, has been
displaced away from the exterior, folded inward, so that the two squares
can now touch directly at the boundary while the displaced material
persists, hidden, along a folded and occluded route that still connects
them exactly as it always did, just no longer where anyone looking at
the finished object can see it.

This can be stated as a general principle, though a modest one: new
visible adjacency does not require new material connectivity. What the
construction changes is which relations among the sheet's points are
expressed at the boundary, not the sheet's underlying connectedness. The
apparent paradox, that two regions never touching in the flat layout end
up touching in the folded object, dissolves once hidden layers are
admitted as a legitimate place for material to be, rather than treated
as an embarrassment to be minimized.

\hypertarget{the-exterior-as-a-quotient}{%
\section{The exterior as a
quotient}\label{the-exterior-as-a-quotient}}

A useful way to state this relationship formally is to treat the
finished, visible object as a coarse quotient of the complete folded
state. The complete folded state includes every layer, every fold, every
point of contact between one part of the sheet and another, all of it.
Most of this detail becomes observationally indistinguishable once only
the exterior is inspected: many different interior arrangements of
hidden layers would produce the same visible exterior, and an observer
confined to looking at the outside of the object has no way to tell
which of these arrangements actually occurred.

The quotient retains exactly what the boundary grammar of Chapter 4
declared mattered: face identity, orientation, and the adjacency
relations among exposed regions. It suppresses everything else, not
because everything else is unimportant to the object's physical
structure, it plainly is important, load-bearing, and literally what
gives the object its bulk, but because the exterior, considered purely
as a visible surface, does not preserve enough information to
reconstruct it.

\hypertarget{surplus-area}{%
\section{Surplus area}\label{surplus-area}}

The preceding two sections established that new exterior adjacency can
be produced without new material connectivity, provided there is
somewhere for the intervening material to go. This section asks how much
material that actually involves, since the answer bears directly on how
demanding a given specification turns out to be in practice.

Define the \textbf{surplus ratio}:

\[\sigma = \frac{\text{Area}(P) - \text{Area}(E_{\text{visible}})}{\text{Area}(E_{\text{visible}})}\]

where \(\text{Area}(P)\) is the total area of the original sheet and
\(\text{Area}(E_{\text{visible}})\) is the combined area of the regions
actually intended to remain exposed. A conventional cube net, built from
close to the minimal area six unit squares require, has a surplus ratio
near zero: essentially all of the sheet becomes visible exterior, give
or take whatever small tabs are needed for joining. The Supercube's
surplus ratio was, by contrast, substantial: the three marked squares
accounted for only a modest fraction of the sheet's total area, with the
remainder, several times larger, constituting pure surplus, none of it
intended for exterior legibility.

A high value of \(\sigma\) is not, by itself, a difficulty. It is a
resource. The surplus material is exactly what makes the vertex-gluing
and adjacency conditions of Chapter 4 achievable in the first place,
since it is the material that gets folded away to create the separation
between the flat layout and the folded exterior that Section 6.2
described. What a high surplus ratio does introduce is a downstream
question: where does all of that material actually go once it has been
folded away, and what does accommodating it cost the rest of the
construction. That question is taken up directly in Section 6.5.

\hypertarget{the-central-tradeoff}{%
\section{The central tradeoff}\label{the-central-tradeoff}}

The material displaced by a high surplus ratio does not vanish; it has
to be accommodated somewhere inside the boundary the exposed faces
define, and accommodating it is not free. This yields the tradeoff at
the center of this chapter's argument:

\[\text{exterior regularity} \;\leftrightarrow\; \text{hidden interior complexity}\]

A cleaner, more regular exterior, flatter faces, sharper edges, more
exact adjacency, does not come for free when the surplus ratio is high.
It is typically purchased by a more complicated interior: a greater
number of hidden layers, sharper curvature concentrated out of view,
more stored elastic energy in the tightly folded material, or a more
elaborate route by which the surplus is threaded through the available
interior space without disturbing the exposed faces it is tucked behind.

This tradeoff is visible directly in the photographic record of the
Supercube's construction. The version of the object in which the three
marked squares present as relatively clean, recognizable faces required
the surrounding material to be folded and crumpled quite densely to fit
inside; the alternative crumpled states, photographed as the same sheet
was carried further into increasingly irregular configurations, show the
same marked regions still present and still identifiable, but no longer
held to the same exterior regularity, with the freed regularity
apparently absorbed as a less constrained, more loosely packed interior.
Neither state is more or less correct relative to the declared
specification. The specification, recall from Chapter 4, said nothing
about the sacrificial material's degree of regularity. Both states are
realizations of the same boundary grammar, purchased at different points
along the same tradeoff.

It is worth naming a related, practical consequence directly:
\textbf{material congestion}. As more surplus is folded inward, the
available interior space fills, and the material already there can begin
to resist further rearrangement, potentially forcing an
already-established face outward, disturbing an adjacency that had
previously been achieved, or blocking a corner from closing correctly.
Section 5.6's discussion of relaxation and repair describes the
operational response to exactly this kind of congestion; this chapter's
contribution is to note that congestion is not an incidental hazard of
infolding but a direct, quantifiable consequence of how much surplus a
given specification's surplus ratio requires the construction to absorb.

\hypertarget{creases-as-a-historical-record}{%
\section{Creases as a historical
record}\label{creases-as-a-historical-record}}

Every fold introduced during the Supercube's construction left a
physical trace on the sheet: a crease, a compressed region of fibers, a
line along which the paper's resistance to bending has been permanently
reduced. These traces persist even where the fold that produced them is
later reversed or relaxed. This section treats them as what they are, a
physical record of the construction's actual history, rather than as
incidental byproducts of reaching the finished shape.

The declared markings, the three squares and their edge and vertex
annotations, state what the finished object was required to achieve. The
accumulated crease network states, after the fact, what path the sheet
was actually taken along in order to achieve it. These are different
kinds of information, and the distinction between them, a declared
target and a witnessed trajectory, is one this monograph will return to
directly in Section 6.8.

A crease is not merely evidence that a particular fold occurred at a
particular moment. It also constrains what can happen afterward. Paper
that has already been creased along one line bends more readily along
that line in the future and more stiffly elsewhere; a region that has
already been compressed into a tight pleat cannot easily be returned to
a flat, unstressed state without visible residual curvature. This means
that the space of configurations reachable from the sheet's current
state depends on the specific history that brought the sheet to that
state, not merely on the sheet's current visible shape. Two sheets that
look identical at some intermediate stage of construction, but that
arrived at that stage by different sequences of folds, may not have the
same set of configurations reachable from that point onward.

This dependency is worth naming explicitly:

\[\text{Reach}(P, h_t)\]

denotes the set of configurations reachable from the sheet's current
state, given that its deformation history up to time \(t\) has been
\(h_t\). The notation makes explicit what the previous paragraph argued
informally: reachability is a property of the pair, sheet and history,
not of the sheet's current shape alone. A construction that has locked a
relation, in the sense of Section 5.5, and creased the surrounding
material accordingly, has narrowed \(\text{Reach}(P, h_t)\) in a way
that a differently-creased sheet at the same superficial stage would not
have narrowed it.

\hypertarget{many-histories-one-boundary}{%
\section{Many histories, one
boundary}\label{many-histories-one-boundary}}

The preceding sections established two facts that combine into this
chapter's central result. First, a given exterior specification can be
satisfied while displacing surplus material along many different routes
through the interior, since the specification says nothing about how
that displacement is to be organized. Second, the specific route
actually taken is recorded, irreversibly to a first approximation, in
the sheet's accumulated crease network, and different routes leave
different, generally distinguishable, records.

Putting these together: many different folding histories, producing many
different interior crease networks and layer arrangements, can realize
the identical exterior boundary specification.

\[\text{many histories} \;\rightarrow\; \text{many interiors} \;\rightarrow\; \text{one boundary class}\]

This is not a hypothetical possibility raised for the sake of argument.
It is directly demonstrated by the Supercube's own photographic record:
the same marked sheet, carried through different further manipulations,
produces both a relatively regular cube-corner exterior and a family of
considerably more irregular crumpled exteriors, all of them consistent
with the same declared face, edge, and vertex conditions, none of them
privileged by the specification over any other. The specification
defines a class of acceptable exteriors, or in the strictest cases a
class of acceptable exact exteriors together with a much larger class of
acceptable approximate ones, and any member of that class, however it
was reached, counts as satisfying it.

\hypertarget{boundary-form-does-not-disclose-interior-history}{%
\section{Boundary form does not disclose interior
history}\label{boundary-form-does-not-disclose-interior-history}}

The consequence of Section 6.7 that matters most for the remainder of
this monograph is epistemological rather than constructional, and it
deserves to be stated as directly as possible.

Given only the finished, folded exterior of an object realized by
constraint-directed infolding, an observer cannot, by inspection of that
exterior alone, recover which of the many possible interior histories
actually produced it. The exterior does not encode this information,
because the specification that governs the exterior was never designed
to encode it; recall from Chapter 4 that the boundary grammar constrains
only exposure, adjacency, orientation, vertex incidence, relegation, and
closure, and says nothing whatsoever about the internal route by which
relegation is achieved. An object satisfying that grammar is, from the
standpoint of the grammar itself, complete and correct regardless of
which member of the many-histories-many-interiors class it happens to
be.

This should be distinguished carefully from a claim that the interior is
unknowable in principle. It is not unknowable; Chapter 8's discussion of
imaging and destructive unfolding takes up exactly how an interior might
be recovered by means other than inspecting the exterior. The claim is
narrower and, for that reason, sharper: the exterior itself, considered
purely as a boundary satisfying a declared grammar, carries no reliable
trace of its own history. Recovering that history requires additional
access, physical or otherwise, beyond what the finished boundary
presents.

This licenses a distinction among several notions of equivalence that
have been used loosely up to this point and that now need to be kept
separate.

Two realizations are \textbf{boundary-equivalent} when they satisfy the
same declared exposure, adjacency, orientation, vertex, and closure
conditions, regardless of anything else about them. This is the weakest
and most important notion for the purposes of this monograph, since it
is the notion the Supercube's specification was actually designed to
secure.

Two realizations are \textbf{structurally equivalent} when they are
boundary-equivalent and additionally share the same pattern of interior
layer contacts, even if the specific creases producing those contacts
differ in detail.

Two realizations are \textbf{crease-equivalent} when they share, in
addition to structural equivalence, the same crease network up to a
specified tolerance, meaning not merely that the same layers touch but
that they touch along recognizably the same fold lines.

Two realizations are \textbf{procedurally equivalent} when they were
produced by the same sequence of operations, in the same order, the
strongest and least useful notion of equivalence for present purposes,
since it is closest to the strong reproducibility Chapter 2 associated
with conventional sequential diagrams.

These four notions form a hierarchy, each one a strictly stronger
requirement than the last, and the Supercube's two documented outcomes,
the relatively regular cube corner and the family of more irregular
crumpled states, are boundary-equivalent to a considerable degree,
agreeing on the exposure, adjacency, and vertex conditions the
specification actually demanded, while being neither structurally,
crease-, nor procedurally equivalent to one another in any strong sense.
This is not a defect in either realization. It is the expected
consequence of specifying at the constraint level, in the sense of
Chapter 2, rather than at the crease or operation level: a
constraint-level specification purchases its sparseness at the price of
exactly this kind of underdetermination, and the underdetermination is
not a loose end to be tied up but the specification doing precisely what
it was built to do.

The chapter's argument can now be closed with the proposition it has
been building toward throughout:

\begin{quote}
Exterior equivalence does not imply structural, crease, or procedural
equivalence.
\end{quote}

A simple, legible exterior can be the visible quotient of a complicated,
irreversible, and largely unrecoverable interior process, and nothing
about the exterior's simplicity licenses an inference about the
simplicity, or even the singularity, of the process that produced it.
Chapter 7 takes up the other side of this same fact, asking what it
means to conduct the search this chapter has been analyzing after the
fact, from the inside, without the benefit of already knowing which
interior history one is going to end up leaving behind.

\part{Embodied Search}

\hypertarget{chapter-7.-folding-without-a-complete-internal-image}{%
\chapter{Folding Without a Complete Internal
Image}\label{chapter-7.-folding-without-a-complete-internal-image}}

\hypertarget{the-search-from-the-inside}{%
\section{The search from the
inside}\label{the-search-from-the-inside}}

Chapter 6 described the space of possible realizations of a boundary
specification from the outside, as an analyst might look back on it once
several outcomes were available for comparison. This chapter asks what
the same search looks like from the inside, to a constructor who has
only the marked sheet, their hands, and the sheet's own resistance to
work with, and who does not yet know, and cannot know in advance, which
member of the many-histories-many-interiors class described in Section
6.7 they are going to end up producing.

This is not a digression from the formal argument. It bears directly on
what kind of specification a constraint-level description, in the sense
of Chapter 2, can actually be expected to support. A specification that
could only be realized by a constructor holding a complete mental image
of the finished object's interior, planned out in advance down to every
crease, would not really be sparse in the sense Chapter 3 claimed for
it; the sparseness would be an illusion, with all the missing detail
simply relocated into the constructor's head rather than genuinely left
open. What this chapter argues is that the Supercube's construction did
not require this, and that its actual method is more interesting, and
more generally useful, than a mental substitute for the missing
specification would have been.

\hypertarget{externalized-visualization}{%
\section{Externalized
visualization}\label{externalized-visualization}}

The marked sheet itself carries the information a constructor needs to
proceed, at any given moment, without requiring that constructor to hold
a complete picture of the finished object in mind. The three squares,
their edge annotations, and their shared vertex label remain physically
present and physically inspectable throughout the construction. A
constructor does not need to remember, in the sense of recalling from
memory, which edge was meant to be the bottom edge; the mark is still
there, on the sheet, wherever the sheet currently is, and can simply be
looked at again.

This matters more than it might first appear to, because it changes what
kind of cognitive demand the construction actually places on whoever is
carrying it out. The demand is not: hold the complete target geometry in
mind and manipulate the sheet to match it. The demand is narrower and
more local: look at the currently visible marks, notice which declared
relations are not yet satisfied, and make whatever move seems likely to
bring one of them closer to being satisfied. The marked sheet functions
as an external store of exactly the information that a complete mental
image would otherwise have had to hold internally.

This point bears specifically, and should be stated carefully, on
aphantasia, the reported absence or marked reduction of voluntary visual
imagery in some individuals. It would be a significant overreach to
suggest that this construction method demonstrates anything about the
spatial reasoning capacities of people who report aphantasia, and no
such claim is made here. What can be said, more modestly, is that a
construction method built around externalized, physically inspectable
marks and iterative local search does not appear to presuppose the kind
of detailed, sustained voluntary visual imagery whose absence aphantasia
describes. Whether or not a given constructor experiences vivid mental
imagery of the finished cube while folding toward it, the marked sheet
supplies, externally, the same information that imagery might otherwise
have supplied internally. This is a claim about what the method
requires, not a claim about what any particular person's imagery does or
does not include, and it should not be read more broadly than that.

\hypertarget{local-solvability}{%
\section{Local solvability}\label{local-solvability}}

A further consequence of working from externalized marks rather than a
complete internal plan is that the construction can proceed by
addressing small, local questions one at a time, rather than requiring
the global configuration to be solved all at once.

At any given moment in the Supercube's construction, a useful next move
could be identified by asking a narrow question: can these two marked
edges be brought closer together without disturbing a relation already
secured? Can the corner that has already been formed be held steady
while the material around it is rearranged? Can a bulge of surplus
material, currently protruding where a marked face needs to sit flat, be
pushed further inward? None of these questions requires an answer to the
larger question of what the finished object will ultimately look like in
its entirety. Each is answerable by inspecting the sheet's current state
against the declared marks and judging, from that comparison alone, what
the next useful move might be.

Global coherence, on this account, is not planned in advance and then
executed. It accumulates, as a consequence of a sequence of locally
useful moves, each one chosen for its immediate contribution to
satisfying some currently unsatisfied relation, none of them requiring a
synoptic view of the whole. This is worth stating plainly because it is
a genuinely different model of how a complex spatial construction can be
carried out than the model implicit in a conventional crease pattern,
which does presuppose that the complete relationship between flat sheet
and folded form is knowable, and known, before folding begins.

\hypertarget{failure-as-information}{%
\section{Failure as information}\label{failure-as-information}}

A fold that does not work, an edge that will not reach its declared
partner, a corner that keeps slipping out of alignment as soon as
attention moves elsewhere, is not, on the account developed here, an
event external to the search process, an interruption to be gotten past
as quickly as possible. It is itself a piece of information about the
current shape of what Section 6.6 called \(\text{Reach}(P, h_t)\): the
set of configurations still reachable given the sheet's history up to
that point.

A blocked fold indicates one of a small number of specific things: that
the surrounding material has become congested, in the sense of Section
6.5, and needs to be redistributed before the fold can proceed; that two
regions have been brought into an orientation genuinely incompatible
with the fold being attempted, requiring one or the other to be
repositioned first; that insufficient slack remains in the surrounding
sheet to permit the motion required, requiring surplus to be drawn in
from elsewhere; or that a relation locked earlier, in the sense of
Section 5.5, is now standing in the way of further progress and needs to
be relaxed, per Section 5.6, before construction can continue. In each
case, the failure is diagnostic. It reveals something true about the
sheet's current reachable set that was not evident before the failed
attempt was made, and it narrows the space of useful next moves rather
than merely wasting the move that failed.

\hypertarget{repair-and-workaround}{%
\section{Repair and workaround}\label{repair-and-workaround}}

The natural response to a blocked fold, on this account, is not to undo
the entire construction back to some earlier, unblocked state, though
that remains available as a last resort. It is to introduce a local
workaround: an additional pleat that provides the slack a
straightforward fold could not; a pocket that temporarily holds a corner
in place while attention moves to a different part of the sheet; a
rotation of an already-folded bundle of material to relieve pressure
against an adjacent, protected face; or a deliberate, bounded
distortion, tolerated for the remainder of the construction, that trades
a small amount of cosmetic regularity for continued progress toward the
relations that actually matter.

These workarounds are local in the specific sense that they enlarge the
reachable configuration space in the immediate vicinity of the
obstruction without requiring the constructor to revisit or reconsider
relations that have already been secured elsewhere on the sheet. This is
what makes the process tractable by hand: a constructor need not hold
the entire history of the construction in active consideration at every
step, only the currently active obstruction and whatever already-secured
relations that obstruction's resolution must not disturb.

\hypertarget{traces-and-self-scaffolding}{%
\section{Traces and
self-scaffolding}\label{traces-and-self-scaffolding}}

Each fold made during construction leaves a trace, in the form of a
crease, an established edge, or a temporary pocket, and these traces are
available to guide subsequent action in a way that goes beyond simply
recording that the fold occurred. An edge that has already been folded
provides a stable reference against which a later fold can be aligned. A
pocket formed to hold one corner in place can, incidentally, provide a
convenient boundary against which an adjacent region can be pressed
flat. The sheet, over the course of construction, progressively builds
its own scaffolding: later moves are made easier, or made possible at
all, by structure the earlier moves happened to leave behind.

This is a form of what is sometimes called stigmergic coordination in
other contexts, where independent agents, or in this case a single agent
acting over time, coordinate not through a shared plan held in common
but through modifications each leaves in a shared environment that the
others, or the agent's own later self, can perceive and respond to. No
single moment of the Supercube's construction contained the complete
solution. The solution, such as it was, emerged from the accumulated
interaction among the declared markings, which never changed, the
sheet's physical resistance, which shaped what could be attempted at
each step, the traces left by earlier folds, which scaffolded later
ones, and the ongoing comparison between the sheet's current state and
the relations still outstanding.

This observation extends naturally, though only briefly needs to be
extended here, to the possibility of multiple constructors, or multiple
robotic manipulators, working on the same sheet without a complete
shared plan. If each agent can perceive the declared marks and the
traces left by prior actions, whether their own or another agent's,
coordination of a genuinely useful kind becomes possible without
requiring the agents to have agreed in advance on a complete sequence of
operations. This possibility is noted here as a natural extension of the
single-constructor account just given, rather than developed as an
independent theory in its own right; the single-constructor case already
contains everything essential to the point this chapter has been making,
and multiplying the number of hands involved does not change the
underlying logic of local, trace-guided, constraint-preserving search.

\hypertarget{what-the-embodied-account-adds-to-the-formal-one}{%
\section{What the embodied account adds to the formal
one}\label{what-the-embodied-account-adds-to-the-formal-one}}

Chapter 5 gave a formal vocabulary, \(\text{INFOLD}\), \(\text{LOCK}\),
\(\text{RELAX}\), \(\text{REPAIR}\), for describing constraint-directed
search as a sequence of operations subject to declared conditions. This
chapter has given the same process an experiential description,
addressed to what it is actually like, and what it actually requires, to
carry that search out by hand.

The two descriptions are not competing accounts of the same phenomenon
so much as two different vantage points on it. The formal vocabulary
specifies what a valid sequence of operations must accomplish, without
saying anything about how a constructor is to discover such a sequence.
This chapter has argued that the discovery process does not require, and
in the Supercube's case did not involve, a complete internal image of
the finished object; it required only a persistent external record of
what mattered, a willingness to proceed by local, diagnostic, reversible
steps, and a sheet whose own accumulated creases progressively
scaffolded the steps that followed.

This has a direct bearing on the practical usefulness of
constraint-directed infolding as a method, beyond its interest as a
piece of formal theory. A specification that can be realized this way,
sparse, externally inspectable, and searchable through local moves
alone, is a specification that does not require its constructor to be an
expert planner capable of anticipating a complete folding sequence in
advance. It requires only the ability to compare a sheet's current state
against a small number of declared marks and to recognize, moment to
moment, whether the gap between them is narrowing. Chapter 8 turns to a
different kind of surface facing a structurally similar problem, one
where no constructor, human or otherwise, is available to carry out this
kind of moment-to-moment comparison at all.

\part{Biological Analogy}

\hypertarget{chapter-8.-cortical-infolding-under-surface-volume-constraint}{%
\chapter{Cortical Infolding Under Surface-Volume
Constraint}\label{chapter-8.-cortical-infolding-under-surface-volume-constraint}}

\hypertarget{a-structural-resemblance-stated-cautiously}{%
\section{A structural resemblance, stated
cautiously}\label{a-structural-resemblance-stated-cautiously}}

This chapter develops an analogy between the paper construction
described in the preceding chapters and the folded structure of the
mammalian cerebral cortex. It is important to state at the outset,
before any of the resemblance is drawn out, exactly what kind of claim
is and is not being made. The claim is structural: both systems involve
a surface whose area substantially exceeds what could be smoothly
accommodated within the volume that contains it, and both resolve this
excess through folding rather than through simple expansion of the
containing boundary. The claim is not mechanistic: this chapter does not
propose, and does not believe, that the specific process by which a
sheet of paper is folded by hand bears any direct resemblance to the
specific biological processes by which a developing cortex actually
folds. Those processes are addressed on their own terms in Section 8.3,
and the differences between them and the paper construction are
addressed directly and at comparable length in Section 8.4. The purpose
of drawing the analogy at all is narrower than a claim of mechanism
would be: it is to illustrate, using a case that can be inspected
directly and folded by hand, a general geometric fact about area-rich
surfaces in volume-limited environments that plausibly bears on the
cortex without being confused with an account of how the cortex actually
achieves it.

\hypertarget{the-shared-problem}{%
\section{The shared problem}\label{the-shared-problem}}

Stated abstractly, the shared problem is this: a surface possesses more
area than the volume in which it is to be contained could accommodate if
that surface remained smooth and unfolded. Something has to give. Either
the containing volume expands to match the surface, or the surface's
exposed extent is reduced by folding it into a smaller effective
envelope, with the folded surface's full area preserved but no longer
laid out flat.

The Supercube's sheet faced this problem directly, and by design: the
three marked squares required a specific exterior configuration, and the
sheet's substantial surplus area, defined formally in Section 6.4, had
nowhere to go except into the interior the marked squares' configuration
made available. The cerebral cortex, over the course of development,
appears to face a structurally comparable problem: the cortical sheet's
surface area increases considerably during a period in which the skull's
interior volume is comparatively constrained, and a smooth, unfolded
cortex of the area actually present in a mature human brain would not
fit within a skull of ordinary size. Folding, in the form of the
familiar pattern of gyri and sulci, the ridges and grooves visible on
the brain's surface, allows substantially more cortical surface to exist
within a given cranial volume than an unfolded sheet of the same area
could.

This much is a shared geometric circumstance, and it is worth being
clear that stating it this way commits to nothing about mechanism. It
states only that both systems face an area-in-excess-of-volume problem
and that both resolve it, whatever the process, by folding rather than
by simple expansion.

\hypertarget{an-expansive-metaphor-and-the-established-mechanistic-accounts}{%
\section{An expansive metaphor, and the established mechanistic
accounts}\label{an-expansive-metaphor-and-the-established-mechanistic-accounts}}

One evocative way to describe a surface's response to a spatial
constraint imposed by its own growth, and by the growth of neighboring
structures around it, borrows the phenomenon sometimes called crown
shyness in forestry: the observation that the crowns of certain tree
species, growing close to one another, tend to stop just short of
touching, leaving a visible channel of open sky between adjacent
canopies. This image is invoked here purely as a metaphor for
constrained, mutually negotiated expansion into a bounded and already
partially occupied space, and it should be understood as nothing more
than that. Crown shyness is not itself an accepted account of cortical
folding, and no claim is made that the mechanisms producing it in forest
canopies bear any relation to the mechanisms producing folds in cortical
tissue. The metaphor's only role is to evoke, vividly, the general idea
of an expanding surface whose form is shaped by the presence of
neighboring constraints rather than by its own growth in isolation, a
general idea the cortex's situation shares even though the specific
dynamics involved are entirely different.

The actual, established explanations for cortical gyrification are
mechanistic and biological, and several distinct candidate mechanisms
currently coexist in the literature without a single one having
displaced the others as a complete account. Differential tangential
expansion proposes that the cortical gray matter expands more rapidly,
or along different constraints, than the white matter beneath it,
producing a mechanical instability comparable to the buckling of a
stiff, expanding layer bonded to a more slowly expanding substrate.
Radial constraint accounts emphasize the geometric confinement imposed
by the shape and growth rate of underlying structures. Axonal tension
hypotheses propose that mechanical tension along developing white-matter
connections pulls strongly interconnected cortical regions toward one
another, contributing to the folding pattern. Additional accounts point
to regional variation in cortical thickness, to the developmental
sequencing and heterogeneity of different cortical areas, to the growth
of underlying white matter volume, and to the geometry of the
ventricular system as further contributing factors. These accounts are
not mutually exclusive, and the current state of the relevant
developmental biology treats gyrification as very likely multicausal,
with different mechanisms contributing at different developmental stages
and in different regions. Adjudicating among them is not the purpose of
this chapter, and no attempt is made here to do so; they are listed to
make clear that a substantial, ongoing, and specifically biological
research program already exists to explain cortical folding, and that
the paper analogy developed in this chapter is not offered as a
contribution to that program or as a substitute for it.

\hypertarget{what-the-analogy-usefully-captures}{%
\section{What the analogy usefully
captures}\label{what-the-analogy-usefully-captures}}

With the mechanistic accounts of Section 8.3 kept firmly separate, the
structural analogy can be stated with some precision.

Both the folded sheet and the folded cortex preserve material continuity
while changing extrinsic spatial relations. Just as Chapter 6 showed
that the Supercube's three marked squares became exterior neighbors
without any new material connection being formed between them, cortical
regions that lie at some geodesic distance from one another along the
intrinsic, unfolded cortical surface can end up in close
three-dimensional proximity once the mature, folded cortex is examined,
without any claim that this proximity reflects a direct developmental or
connective relationship between those regions. Conversely, and just as
significantly, two points that are close together on the folded
exterior, sitting on either side of a sulcus, for instance, may be quite
distant from one another along the cortical sheet's own intrinsic
surface. Three-dimensional spatial proximity in a folded structure,
whether paper or cortex, is not a reliable guide to intrinsic adjacency,
and intrinsic adjacency is not a reliable guide to three-dimensional
proximity.

This licenses a distinction, already familiar from the paper case, among
several different notions of nearness that a cortical structure's
geometry can support simultaneously and that do not, in general,
coincide: material or intrinsic surface adjacency, three-dimensional
geometric proximity after folding, connective or axonal proximity,
developmental or lineage-based proximity, and functional or task-based
association. The paper construction demonstrates cleanly, because it can
be unfolded and directly inspected, that these several notions of
nearness genuinely can diverge from one another once a surface has been
folded; it offers this as a general geometric possibility available to
any folded surface, cortex included, without asserting anything further
about which of these notions of nearness the cortex's actual folding
pattern happens to preserve or disrupt in any specific instance.

A further point of resemblance concerns history. Section 6.6 argued that
a folded sheet's crease network functions as a partial, physically
embedded record of its own folding history, and that this record
constrains, without fully determining, what future configurations remain
reachable. A mature, folded cortex likewise carries the residue of its
own developmental history in its specific pattern of gyri and sulci, a
pattern shaped by the developmental constraints, whichever combination
of the mechanisms described in Section 8.3 turns out to have been
operative, that were actually in force during its formation. And just as
Section 6.7 showed that the same declared boundary specification could
be reached by many different interior histories, it is a live
possibility, consistent with current developmental biology though not
established as fact by anything in this chapter, that similar mature
folding patterns could arise from somewhat different combinations of the
contributing developmental parameters, and that small differences
arising early in development could be amplified into more substantial
differences in the mature folding pattern. This is offered as a
structural possibility the paper analogy makes vivid, not as an
empirical claim about cortical development that this monograph is in a
position to establish.

\hypertarget{where-the-analogy-breaks}{%
\section{Where the analogy breaks}\label{where-the-analogy-breaks}}

Several differences between the two systems are fundamental enough that
the analogy should not be pressed past the structural point Section 8.4
draws from it, and each is worth stating plainly rather than left as an
implicit qualification.

Paper is a passive, manufactured material, externally manipulated by an
agent, the constructor, who exists outside the sheet and who folds it
according to a plan, however sparse. Cortical tissue is living matter
that grows, differentiates, remodels itself, develops its own
vasculature, and interacts mechanically with the tissues surrounding it,
all without an external agent folding it according to any plan at all;
whatever organizing structure the process has is generated from within
the developing tissue and its immediate environment, not imposed from
outside by a hand holding a target specification.

The Supercube began with an explicit symbolic specification: three
marked squares, each carrying declared edge roles and a declared shared
vertex, fixed in advance of any folding and legible throughout the
construction as an external reference the constructor could return to at
any point. Cortical development does not appear to require, and there is
no evidence that it involves, any comparable symbolic representation of
the mature cortex's eventual folded form, held somewhere in advance of
the folding process and consulted during it. Whatever determines the
mature folding pattern is encoded, if the word encoded is even
appropriate, in genetic, molecular, and mechanical factors operating
locally and dynamically during development, not in anything resembling
the three discrete, externally legible marks that organized the paper
construction.

Uncontrolled crumpling, of the kind demonstrated in the alternative,
more irregular states the Supercube's sheet was carried into in Section
6.5's discussion, can produce sharp ridges, localized regions of high
curvature, and structural singularities that a passive sheet of paper
tolerates without consequence but that living tissue plausibly could
not, given constraints on cell viability, blood supply, and mechanical
stress that a folded sheet of paper simply does not face. Cortical
folding, whatever its precise mechanism, evidently proceeds under
biological constraints on curvature, thickness, and structural integrity
that have no counterpart in the paper case, and there is no reason to
expect the geometric strategies available to an unconstrained sheet of
paper to all be available to living cortical tissue.

\hypertarget{a-restrained-conclusion}{%
\section{A restrained conclusion}\label{a-restrained-conclusion}}

Given the qualifications of Section 8.5, the appropriate claim this
chapter is entitled to make is considerably narrower than a reader might
have expected on first encountering the analogy, and it is worth stating
in a form precise enough to prevent the analogy from being carried
further than this chapter intends.

\begin{quote}
The Crinkle-Cut Supercube is not a model of the mechanism of cortical
gyrification. It is a model of the more general geometric fact that
infolding allows an area-rich surface to inhabit a volume-limited
environment while generating spatial relations absent from its unfolded
presentation.
\end{quote}

Everything specific to how the cortex actually achieves this, which of
the mechanisms surveyed in Section 8.3 predominates, in what
combination, at what developmental stage, in what region, is a question
for developmental biology to answer, and nothing in this chapter
contributes to answering it. What the paper construction contributes is
a directly inspectable, foldable, unfoldable demonstration that the
general strategy, preserving material continuity while displacing
surplus surface away from a constrained exterior, is geometrically
available at all, and that a folded structure's exterior configuration,
once achieved, will not by itself disclose the specific history or
mechanism that produced it. This last point is the one thread that
connects this chapter most directly back to the argument of Chapter 6,
and it is the thread Chapter 9 now picks up in more general form, asking
what can be said, formally, about when a target exterior configuration
of this kind is achievable at all.

\part{General Principle and Consequences}

\hypertarget{chapter-9.-continuation-reachability-and-the-crinkle-cut-principle}{%
\chapter{Continuation, Reachability, and the Crinkle-Cut
Principle}\label{chapter-9.-continuation-reachability-and-the-crinkle-cut-principle}}

\hypertarget{admissibility-at-each-stage}{%
\section{Admissibility at each
stage}\label{admissibility-at-each-stage}}

At any moment during the construction of the Supercube, the sheet
occupied some particular partially folded state, and from that state,
some further moves were available and others were not. This is true
trivially, in the sense that any physical process unfolds through a
sequence of intermediate states, but it is worth making the observation
precise, because the distinction between moves that remain available and
moves that do not is what the rest of this chapter is built on.

Call a move \textbf{locally admissible} if making it does not
irreparably violate a protected constraint, in the sense of Section 5.4,
and does not damage the material integrity of the sheet itself, through
tearing or an impermissible crease. A locally admissible move need not
bring the construction closer to completion. It only needs to avoid
closing off possibilities that were previously open.

\hypertarget{local-admissibility-is-not-enough}{%
\section{Local admissibility is not
enough}\label{local-admissibility-is-not-enough}}

It is tempting to suppose that a construction proceeds successfully
whenever each individual move, taken on its own, is locally admissible:
if no single step damages anything already secured, the construction as
a whole ought to succeed. This is false, and it is worth being precise
about why, since the failure mode involved is exactly the kind of
failure the Supercube's own construction had to navigate around.

A move can be locally admissible, in the narrow sense of Section 9.1,
while still closing off the global path to completion. Bringing two of
the three marked faces into close alignment, for instance, does not
itself violate any protected constraint and may look, examined in
isolation, like straightforward progress. But if that alignment traps a
mass of surplus material in a position from which the third marked face
cannot subsequently be brought into its required adjacency without
either disturbing the first two or exceeding the sheet's capacity to
bend without tearing, then the move was locally admissible and globally
harmful at once. Nothing about examining the move on its own terms, at
the moment it is made, reveals this consequence; the consequence only
becomes visible once later moves are attempted and found to be blocked.

This distinction between local admissibility and \textbf{global
reachability}, the continued existence of at least one path from the
current state to a configuration satisfying every declared constraint,
is the central technical idea of this chapter. A construction can
proceed for a considerable time making only locally admissible moves and
still arrive at a state from which no completion is reachable at all.
Recognizing this possibility, and organizing the construction to guard
against it, is a large part of what made the Supercube's own
construction, as reconstructed in Chapter 1, proceed by establishing the
vertex-gluing condition early and treating it as a stable anchor, in the
sense of Section 5.5's \(\text{LOCK}\) operation, rather than by
pursuing whichever local improvement happened to present itself first.

\hypertarget{a-hierarchy-of-protected-distinctions}{%
\section{A hierarchy of protected
distinctions}\label{a-hierarchy-of-protected-distinctions}}

Because not every locally admissible move preserves global reachability,
and because a construction carried out by hand cannot exhaustively
verify global reachability before every move, some principle is needed
for deciding which relations to protect most strongly when a choice has
to be made between preserving one relation and preserving another. The
Supercube's construction, examined in retrospect, appears to have
followed an implicit ranking of exactly this kind, and it is useful to
state that ranking explicitly:

\begin{gather*}
\text{material continuity} \;>\; \text{selected-face identity} \;>\; \text{required incidence} \\
{}\;>\; \text{edge orientation} \;>\; \text{planarity} \;>\; \text{cosmetic regularity}
\end{gather*}

Read from left to right, this hierarchy states that the sheet's basic
material continuity, its status as a single connected surface without
tears, is protected above everything else, since its loss would
undermine the entire theoretical apparatus of Chapter 6, which depends
on new exterior adjacency being achievable without new material
connection. Below that, the identity of the three selected faces, which
regions of the sheet count as the declared exterior at all, is protected
more strongly than any specific relation among them. Below that again,
the required incidence at the shared vertex and the required edge
orientations are protected in turn, followed by the planarity of the
individual faces, and finally by the cosmetic regularity of the finished
object's overall appearance, which sits at the bottom of the hierarchy
and was, in the Supercube's actual case, the distinction most freely
sacrificed.

This hierarchy explains a fact that would otherwise be puzzling: why the
finished Supercube, visibly rough and imperfect, nonetheless counts as a
successful realization of its declared specification. Exact planarity,
the fifth-ranked distinction, was sacrificed; cosmetic regularity, the
sixth, was sacrificed considerably further. But material continuity,
face identity, vertex incidence, and edge orientation, the four
higher-ranked distinctions, were all preserved. Measured against the
hierarchy rather than against an implicit standard of overall neatness,
the construction achieved exactly what it was organized to achieve, and
its visible roughness is the direct, expected consequence of having
correctly prioritized higher-ranked distinctions over lower-ranked ones
when a choice between them arose, rather than evidence of an unfinished
or poorly executed construction.

\hypertarget{repair-as-renewed-reachability}{%
\section{Repair as renewed
reachability}\label{repair-as-renewed-reachability}}

Section 5.6 introduced \(\text{RELAX}\) and \(\text{REPAIR}\) as
operations permitting a locked relation to be temporarily loosened and
subsequently restored. This chapter's distinction between local
admissibility and global reachability gives those operations a sharper
justification than mere convenience.

A repair succeeds, in the sense relevant here, not merely when it
restores the relation's prior appearance, but when it reopens a path
toward eventual completion that had been at risk of closing, while not
sacrificing any distinction ranked above the one being repaired in the
hierarchy of Section 9.3. A repair that restored a relation's cosmetic
appearance while permanently foreclosing the construction's ability to
satisfy a higher-ranked distinction elsewhere would not be a successful
repair in this sense, whatever it looked like at the moment it was made.
Repair, understood this way, is not restoration for its own sake. It is
a targeted reopening of global reachability, undertaken at the specific
point where reachability had become locally threatened, and evaluated by
whether the threat is actually resolved rather than by whether the
repaired region looks as it did before.

\hypertarget{the-finished-object-as-a-certificate}{%
\section{The finished object as a
certificate}\label{the-finished-object-as-a-certificate}}

Once a construction reaches a state satisfying every constraint in its
declared boundary grammar, that finished object stands as more than a
completed artifact. It is a physical certificate that the declared
specification was, in fact, satisfiable, given the sheet, the material,
and the permitted operations actually available.

This is worth stating with some formality, because the distinction
between demonstrating that a specification is satisfiable and merely
conjecturing that it might be is exactly the distinction Chapter 1 drew,
in its own terms, between a constructive existence claim and a general
theorem. The Supercube's finished form does not certify that every
non-contiguous three-face assignment is realizable; it certifies only
that this one was, by the direct and irrefutable method of having
actually been realized. The hidden folds inside the object, invisible
from the exterior in exactly the sense Chapter 6 described, are the
material trace of the successful search: a physical record, though not,
per Section 6.8, a legible one, that some path through the space of
admissible deformations was actually found and actually taken.

\hypertarget{a-formal-statement-the-crinkle-cut-principle}{%
\section{A formal statement: the Crinkle-Cut
Principle}\label{a-formal-statement-the-crinkle-cut-principle}}

The preceding sections license a formal statement of what the
Supercube's construction actually demonstrates, phrased with the
deliberate restraint the word ``some'' is doing real work to provide:

\begin{quote}
\textbf{The Crinkle-Cut Principle.} For some connected sheets containing
spatially separated, locally framed regions, a target exterior adjacency
among those regions can be realized without cutting the sheet, when the
complementary region is permitted sufficient infolding, overlap,
curvature, and hidden interior occupation.
\end{quote}

The word ``some'' in this statement is not a placeholder for a stronger
claim the monograph is unable to fully justify. It is the precise scope
the constructive evidence actually supports. One sheet, with one
specific assignment of three separated, framed regions, was successfully
realized. This licenses the existence claim above, for that class of
sheets and assignments broadly similar to the one demonstrated, and it
licenses nothing stronger without further argument or further
construction.

\hypertarget{a-stronger-conjecture}{%
\section{A stronger conjecture}\label{a-stronger-conjecture}}

A more ambitious claim can nonetheless be stated, clearly marked as a
conjecture rather than as a result this monograph has established:

\begin{quote}
\textbf{Conjecture.} Broad families of finite boundary specifications
are realizable from sufficiently large connected sheets, subject to
constraints imposed by orientation, thickness, material continuity, edge
length, self-contact, interior capacity, and tolerance.
\end{quote}

This conjecture generalizes the Crinkle-Cut Principle from the
demonstrated case of three separated square regions to a much wider
class of possible boundary specifications, arbitrary numbers of regions,
arbitrary shapes, arbitrary target adjacency graphs, while making
explicit that realizability cannot be expected to hold unconditionally.
The qualifying clause is not a formality. Section 9.8 takes up, in turn,
several of the specific ways a boundary specification can fail to be
realizable regardless of how much surplus surface and permitted
infolding is made available, and these obstructions are exactly what the
conjecture's qualifying clause is intended to acknowledge in advance.

\hypertarget{obstructions}{%
\section{Obstructions}\label{obstructions}}

Several distinct failure modes deserve to be named individually, since
they arise for different reasons and would need to be ruled out
separately in any attempt to establish the conjecture of Section 9.7 for
a specific class of specifications.

A specification may demand \textbf{contradictory edge roles}: two edges
each required, by the declared orientation conditions, to occupy roles
that cannot simultaneously be satisfied by any consistent assignment of
orientation to the intervening sheet. No amount of surplus material or
permitted infolding resolves a contradiction of this kind, since the
obstruction is combinatorial rather than geometric.

A specification may demand adjacency among regions separated by
\textbf{insufficient connecting area}, requiring the intervening
material to be compressed or displaced beyond what its actual area
permits, regardless of how it is folded. This obstruction is geometric
rather than combinatorial, and it interacts directly with the surplus
ratio \(\sigma\) defined in Section 6.4: a specification with too low a
surplus ratio, relative to what the declared adjacency and vertex
conditions actually demand, may be unrealizable for reasons of raw
material insufficiency alone.

A specification may be obstructed by \textbf{excessive material
thickness}, where the sheet's actual physical thickness prevents the
tight folding radii a given configuration requires, an obstruction with
no counterpart in the idealized zero-thickness surface of Section 3.1
but a very real one for any physical sheet.

A specification may demand \textbf{impossible orientation requirements},
asking a region to be simultaneously reflected and unreflected relative
to some other region, in a way no continuous deformation of a connected
sheet can satisfy without tearing.

A specification may run into \textbf{inadequate interior volume}, where
the target exterior's enclosed volume is simply too small to contain the
surplus material the specification's surplus ratio requires to be
relegated, an obstruction distinct from insufficient connecting area,
since it concerns not whether the material can reach the interior but
whether the interior can hold it once it arrives, closely related to the
congestion discussed in Section 6.5.

A specification may produce a \textbf{self-locking trajectory}, in which
every locally admissible sequence of moves available from some
intermediate state leads to a configuration from which no further
locally admissible move preserves global reachability, an obstruction
that depends not on the specification alone but on the specific sequence
of moves attempted, and that a differently-ordered construction,
following Section 9.3's hierarchy more carefully, might have avoided
even where a less careful one could not.

Finally, a specification may simply demand \textbf{geometric exactness
incompatible with crumpling}, requiring a degree of precision in the
final exterior that the permitted family of operations, which after all
includes buckling and crumpling among its admissible deformations,
cannot reliably deliver, even where a less exacting tolerance would have
been readily achievable by the same means.

\hypertarget{ideal-mechanical-and-practical-realizability}{%
\section{Ideal, mechanical, and practical
realizability}\label{ideal-mechanical-and-practical-realizability}}

It is worth closing this chapter by distinguishing three distinct senses
in which a specification can be said to be realizable, since the
obstructions of Section 9.8 do not all apply equally to each.

A specification is \textbf{ideally realizable} if some continuous
deformation of an idealized, zero-thickness, infinitely flexible surface
satisfies it, without regard to whether any physical material could
actually achieve that deformation. A specification is
\textbf{mechanically realizable} if some physically plausible
deformation of a sheet with realistic thickness, stiffness, and plastic
behavior satisfies it, accounting for the obstructions of Section 9.8
that arise specifically from material properties. A specification is
\textbf{practically realizable} if a construction process, whether
carried out by hand, by a computational search of the kind Chapter 10
briefly considers, or by some other method, can actually be found to
locate a satisfying deformation within a reasonable amount of effort,
distinct from the mere fact that such a deformation exists somewhere in
the space of possibilities.

These three notions can, and in general do, come apart. A specification
might be ideally realizable while being mechanically unrealizable for
any sheet of nonzero thickness; it might be mechanically realizable
while remaining practically elusive, its satisfying deformation existing
but proving difficult for any actual search process, human or
computational, to locate. The Supercube's own specification was
demonstrated to be practically realizable, by the strongest possible
evidence such a claim admits, an actual, physical, hand-executed
construction that reached it. Nothing in this chapter should be read as
extending that demonstrated practical realizability to any specification
beyond the one actually constructed, and the Crinkle-Cut Principle of
Section 9.6 is worded precisely to reflect this limit.

Chapter 10 closes the monograph's substantive argument by briefly
indicating how the questions this chapter has posed, but deliberately
not answered in general form, might be pursued further: through
systematic variation of the specification, through computational search
over the space of admissible deformations, and through comparison of
independently constructed realizations of the same declared boundary
grammar.

\hypertarget{chapter-10.-further-work}{%
\chapter{Further Work}\label{chapter-10.-further-work}}

\hypertarget{what-remains-open}{%
\section{What remains open}\label{what-remains-open}}

The preceding chapters have established a constructive existence claim,
a formal vocabulary adequate to state it precisely, a theoretical
account of why a satisfied boundary specification cannot disclose the
interior history that satisfied it, and a hedged principle, together
with an explicit list of the ways that principle's stronger conjectured
form could fail to hold. What they have not done, and do not attempt to
do, is settle any of the questions that principle and conjecture leave
open. This short chapter indicates, briefly and without developing any
of them into a research program in their own right, the directions in
which that settling might proceed.

\hypertarget{systematic-variation}{%
\section{Systematic variation}\label{systematic-variation}}

The most direct way to test how far the Crinkle-Cut Principle of Section
9.6 extends beyond the single demonstrated case is to vary the
specification systematically and observe where realizability holds and
where it fails. This means constructing further sheets with different
numbers of selected regions, different relative placements and
separations, different assigned edge orientations, and different surplus
ratios, and attempting, by the same hand-folding methods described in
Chapter 7, to realize each one. Success and failure should both be
recorded, including the specific obstruction encountered in each failed
attempt, drawn from the list in Section 9.8, since a pattern of which
obstructions recur under which conditions would itself be informative
about where the boundary of practical realizability actually lies.

A closely related line of variation concerns material rather than
specification: repeating a fixed specification across sheets of
differing thickness, stiffness, and plastic behavior, to determine how
much of the boundary between success and failure is a property of the
abstract specification and how much is a property of the particular
material used to realize it. Section 9.9's distinction between ideal,
mechanical, and practical realizability makes exactly this kind of
variation necessary to disentangle, since a specification's mechanical
realizability depends on material properties in a way its ideal
realizability does not.

\hypertarget{independent-realizations-of-the-same-specification}{%
\section{Independent realizations of the same
specification}\label{independent-realizations-of-the-same-specification}}

Chapter 6 argued, from a single sheet carried through two documented
outcomes, that the same boundary specification can be satisfied by more
than one interior history. A stronger version of this claim would ask
several independent constructors, given the same specification and no
knowledge of one another's methods, to realize it separately, and would
then compare the resulting interiors, by the destructive unfolding or
sectioning such a comparison would require, to determine how much
genuine diversity of interior structure a single boundary class actually
admits in practice, as opposed to in principle. This would give
empirical content to the equivalence hierarchy introduced in Section
6.8, boundary, structural, crease, and procedural equivalence, by
measuring how far apart independently produced realizations of the same
specification actually tend to fall along that hierarchy.

\hypertarget{computational-and-robotic-search}{%
\section{Computational and robotic
search}\label{computational-and-robotic-search}}

The \(\text{INFOLD}\) operator of Chapter 5 was framed throughout as a
search instruction rather than a determinate procedure, and nothing in
principle restricts that search to a human hand. A computational model
could represent the selected regions as rigid or near-rigid panels
connected by a deformable sheet, and search, by whatever optimization
method proves suitable, for a configuration satisfying a declared
boundary grammar while minimizing some combination of bending energy,
self-collision, and interior congestion. Whether such a search reliably
locates realizations that a human constructor would also find, or
instead discovers structurally different solutions unavailable to
hand-folding, would itself bear on the practical-realizability question
raised in Section 9.9.

A robotic implementation, capable of perceiving the declared marks on a
physical sheet and manipulating it directly, would test a further and
more demanding question: whether the constraint-level notation developed
in Chapter 4 carries enough information to guide a search process that
has no access to the embodied, trace-guided intuitions described in
Chapter 7, or whether some of what made the Supercube's construction
tractable by hand depended on forms of tactile feedback and local
judgment that a notation, however carefully specified, cannot fully
substitute for.

\hypertarget{a-deliberate-limit-on-this-chapters-scope}{%
\section{A deliberate limit on this chapter's
scope}\label{a-deliberate-limit-on-this-chapters-scope}}

None of the preceding directions is developed here beyond the brief
indication given above. Detailed experimental templates, replication
protocols, measurement instruments, and standardized recording forms
belong to a separate, purpose-built replication supplement rather than
to the body of this monograph, should such a supplement later be
prepared. Including that material here would extend the monograph well
past the scope this final part is meant to occupy, and would risk
obscuring, under a mass of procedural detail, the more limited claim the
monograph has actually been concerned to establish throughout: that one
nontrivial constraint-directed infolding was successfully demonstrated,
and that its formal and theoretical consequences, worked out across the
preceding chapters, hold independently of whether any of the further
work sketched in this chapter is ever carried out at all.

\hypertarget{conclusion.-the-visible-solution}{%
\chapter*{Conclusion: The Visible Solution}
\addcontentsline{toc}{chapter}{Conclusion: The Visible Solution}\label{conclusion.-the-visible-solution}}

\hypertarget{returning-to-the-object}{%
\section{Returning to the object}\label{returning-to-the-object}}

It is worth returning, at the end of this monograph, to the object
itself, rather than to the formal apparatus that has accumulated around
it across the preceding chapters. Three squares were marked on a sheet
considerably larger than the three of them combined. Each square carried
a small amount of annotation: which edge was to be a bottom edge, which
a side edge, which corner was to meet the corresponding corners of the
other two squares at a single shared vertex. Nothing else about the
sheet was specified in advance.

Nothing about the shape of the finished object was declared, in the
sense that a cube net declares it, by laying out beforehand which
regions must touch which. What was declared were fragments: three local
frames and the relations required to hold among them. The rest of the
sheet, unmarked and unaddressed by anything in the specification, was
left to become whatever it needed to become in order for those fragments
to be brought into their required configuration.

\hypertarget{what-the-construction-preserved-and-what-it-surrendered}{%
\section{What the construction preserved and what it
surrendered}\label{what-the-construction-preserved-and-what-it-surrendered}}

The finished object honored the fragments. The three marked squares
ended up mutually adjacent, correctly oriented, and incident to a shared
corner, exactly as their annotations had specified. Everything else
about the object, the exact route the surplus material took to reach the
interior, the specific creases that route left behind, the object's
overall neatness or lack of it, was never specified and was never, in
any meaningful sense, decided in advance. It was discovered, by hand,
through the kind of local, trace-guided search Chapter 7 described, and
it could have been discovered differently, by a different sequence of
moves, without the finished object thereby failing to satisfy the
specification that governed it.

This is the sense in which the object's visible roughness is not a
shortfall to be apologized for. Measured against the hierarchy of
protected distinctions set out in Chapter 9, the construction preserved
everything ranked above cosmetic regularity and sacrificed cosmetic
regularity itself, which was, from the specification's own standpoint,
exactly the correct trade to make. A cleaner object would not have been
a more successful realization of this specification. It would have been
evidence that a different, more demanding specification had been in
force, one this monograph never claimed to be testing.

\hypertarget{what-generalizes-and-what-does-not}{%
\section{What generalizes and what does
not}\label{what-generalizes-and-what-does-not}}

Sparse boundary constraints, this monograph has argued, can govern a
physical transformation without determining its trajectory. This is not
a paradox, once the distinction between a declared relation and a
realized history, developed formally in Chapter 6 and experienced
directly in Chapter 7, is kept in view. A specification that declares
only a small number of relations leaves correspondingly many degrees of
freedom open, and a physical sheet, permitted to fold, pleat, and
crumple freely within those open degrees of freedom, will find some
resolution of them, though not a predetermined one and not necessarily
the same one twice.

The monograph has been careful, throughout, to keep this claim at the
scale the evidence actually supports. One sheet, with one declared
specification, was successfully realized, more than once, in more than
one interior configuration, all of them satisfying the same declared
exterior relations. This licenses the Crinkle-Cut Principle of Chapter 9
in its stated, deliberately narrow form, and it licenses the stronger
conjecture stated alongside it only as a conjecture, not as a result
this monograph has established. The cortical analogy of Chapter 8 was
developed under the same discipline: a structural resemblance worth
drawing out precisely because it is structural, and precisely because it
was kept apart, explicitly and at comparable length, from any claim
about mechanism that the paper construction cannot support and was never
meant to.

\hypertarget{the-hidden-disorder-is-not-the-opposite-of-the-visible-relation}{%
\section{The hidden disorder is not the opposite of the visible
relation}\label{the-hidden-disorder-is-not-the-opposite-of-the-visible-relation}}

The object's visible simplicity, three legible faces meeting correctly
at a corner, and its hidden disorder, a densely folded and largely
unrecoverable mass of surplus material, are not opposed to one another,
and it would misread the entire argument of this monograph to treat them
as though they were. The hidden disorder is not an unfortunate cost paid
for the visible relation, tolerated because no better alternative was
available. It is the specific mechanism by which the visible relation
became achievable at all. A specification this sparse, governing regions
this widely separated, could not have been realized by a sheet forbidden
from crumpling; the freedom to crumple, and the resulting interior
disorder that freedom produces, is not a defect the construction had to
work around but the very capacity the construction depended on
throughout.

This is the point at which the monograph's central formulation should be
restated, now that the chapters supporting it have been laid out in
full:

\begin{quote}
The boundary is not the whole object and does not contain its complete
explanation. It is the visible residue of constraints successfully
preserved through a larger, partly hidden history of deformation.
\end{quote}

\hypertarget{a-closing-note-on-what-has-been-offered}{%
\section{A closing note on what has been
offered}\label{a-closing-note-on-what-has-been-offered}}

This monograph has offered a description, a notation, a theoretical
account of what a satisfied boundary specification can and cannot
disclose about its own history, a restrained biological analogy, and a
formally hedged principle, together with an explicit account of where
that principle's stronger form might fail. It has not offered a complete
theory of paper mechanics, a general algorithm for realizing arbitrary
boundary specifications, or a mechanistic account of cortical
development, and it has said so plainly at every point where the
temptation to claim more than this arose.

What it has offered, at bottom, is an argument built from a single
physical object: that a boundary can be specified more sparsely than a
complete crease pattern requires, that the gap between a sparse
specification and a finished physical form can be closed by permitting
unconstrained material to absorb whatever the specification leaves open,
and that the resulting object, whatever else can be said about it,
stands as a record, though not a legible one, of exactly this having
been done. The three marked squares are still there, on the finished
object, wherever it ends up. Everything else about it was found, not
planned, and the finding is the part of the construction this monograph
has tried, across its length, to take seriously as a subject in its own
right.

\appendix

\hypertarget{appendix-a.-infolding-symbol-set}{%
\chapter{Infolding Symbol
Set}\label{appendix-a.-infolding-symbol-set}}

This appendix collects, in one place, the graphic notation introduced in
Section 4.9. It is intended as a working reference for marking a
physical sheet or a diagram of one, not as a restatement of the
reasoning behind each symbol, which appears in the body of Chapter 4.

\hypertarget{protected-faces}{%
\section{Protected faces}\label{protected-faces}}

A selected region \(E_i\) is drawn with a heavy outline, distinguishing
it at a glance from unmarked, free material. Each protected face is
given a short identifier, \(F_1\), \(F_2\), \(F_3\), and so on, placed
inside or adjacent to the outlined region.

\hypertarget{free-material}{%
\section{Free material}\label{free-material}}

Regions belonging to the infoldable remainder \(I\) are left unmarked,
or given only a light, uniform shading if some visual indication of
extent is useful. The absence of a heavy outline is itself the operative
signal: it declares that the region's internal crease structure lies
outside the specification, not merely that it has not yet been drawn in.

\hypertarget{edge-role-ticks}{%
\section{Edge-role ticks}\label{edge-role-ticks}}

Edge roles are marked near the relevant edge, close to the outline of a
protected face, using short ticks rather than the mountain- or
valley-fold lines of conventional notation, which already carry a
different meaning.

\begin{itemize}
\tightlist
\item
  A single solid tick marks an edge assigned the bottom role, \(B\).
\item
  A double tick marks an edge assigned the side role, \(S\).
\item
  A triple tick marks an edge assigned the top role, \(T\), where a
  construction requires one.
\item
  A matching pair of small letters, placed near two edges belonging to
  different protected faces, marks an adjacency condition
  \(A(E_i{:}e_a,\, E_j{:}e_b)\) between them: the two edges carrying the
  same letter are the two required to meet.
\item
  A small \(\times\) near an edge marks that edge as prohibited from
  exterior exposure, the role \(X\).
\end{itemize}

\hypertarget{vertex-labels}{%
\section{Vertex labels}\label{vertex-labels}}

A small circled label, such as \(V_1\), placed at a corner of a
protected face, indicates that this corner participates in a
vertex-gluing condition. The same label, repeated at a corner of each of
the other protected faces required to meet it, indicates that all
corners carrying that label must become incident to one shared exterior
vertex in the realized object. An unrepeated label, appearing at only
one corner, carries no such requirement and should not be used; every
vertex label in a complete diagram should appear at two or more corners.

\hypertarget{exposure-status}{%
\section{Exposure status}\label{exposure-status}}

Where a region's exposure status is not simply assumed from its being
marked as protected or left free, a short label may be attached directly
to the region:

\begin{itemize}
\tightlist
\item
  \(\text{OUT}\), for a region intended to be principally exterior.
\item
  \(\text{IN}\), for a region intended to be principally interior.
\item
  \(\text{FREE}\), for a region where either status is acceptable to the
  specification.
\item
  \(\text{SEAM}\), for a region permitted to appear only as a narrow
  boundary strip or overlap, rather than as a full exterior face.
\end{itemize}

\hypertarget{phase-annotations}{%
\section{Phase annotations}\label{phase-annotations}}

Where a construction is more easily communicated in stages, brief phase
labels may be written separately from the geometric marks, typically in
the margin or on an accompanying line of text rather than on the sheet
itself: for example, \emph{Phase 1: realize \(V_1\)}, followed by
\emph{Phase 2: align the bottom-edge cycle}, followed by \emph{Phase 3:
relegate the remaining surplus}. These annotations record the order in
which constraints are intended to be secured. They do not record, and
should not be read as recording, the exact hand motions by which any one
phase is to be carried out.

\hypertarget{compatibility-note}{%
\section{Compatibility note}\label{compatibility-note}}

None of the symbols above are intended to replace the standard
mountain-fold, valley-fold, reversal, sink, and rotation symbols of
conventional origami diagramming. A single diagram may combine both
vocabularies freely: conventional symbols specifying a particular local
operation, and the symbols above specifying what that operation is being
performed in order to achieve or preserve. Section 4.10 discusses this
compatibility at greater length.

\hypertarget{appendix-b.-formal-definitions}{%
\chapter{Formal
Definitions}\label{appendix-b.-formal-definitions}}

This appendix gathers, without re-deriving, the formal terms introduced
across Chapters 3 through 6 and Chapter 9, for reference. Each entry
cites the section in which the term is originally motivated and
explained.

\textbf{Substrate (\(P\)).} The connected material sheet on which a
construction is carried out, treated for formal purposes as a connected,
bounded two-dimensional surface initially lying flat in the plane.
(Section 3.1)

\textbf{Selected region (\(E_i\)).} One of a finite collection
\(E = \{E_1, \ldots, E_n\}\) of regions marked in advance on the
substrate, intended to remain legible on the finished object's exterior.
(Section 3.2)

\textbf{Infoldable complement (\(I\)).} The set-theoretic complement
\(I = P \setminus E\): the portion of the substrate not selected for
exterior legibility, whose final geometry is not declared in advance.
Not to be understood as disposable or unimportant material. (Section
3.2)

\textbf{Local frame (\(F_i\)).} A pair \(F_i = (E_i, o_i)\), where
\(o_i\) is orientation data assigning target roles to some or all of the
boundary features of \(E_i\), independent of the region's current
position or its distance from other selected regions. (Section 3.3)

\textbf{Edge role.} A symbolic label attached to an edge of a local
frame, drawn from a vocabulary including \(B\) (bottom), \(S\) (side),
\(T\) (top), \(J_k\) (joining edge, indexed), \(V_k\) (vertex-incident
edge, indexed), and \(X\) (exposure-prohibited). (Section 3.5)

\textbf{Partial specification.} A local frame or boundary grammar in
which some edges or relations are deliberately left unlabeled,
constituting a positive declaration that the corresponding feature's
final configuration is not being constrained, rather than an omission.
(Section 3.5, Section 3.6)

\textbf{Realization (\(f\)).} A deformation carrying the marked
substrate \(P\) into a folded state, evaluated against whether it
satisfies a declared constraint set \(C\). (Section 4.1)

\textbf{Boundary grammar.} The set of constraint types, exposure,
adjacency, vertex-gluing, orientation, relegation, and closure, taken
together, by which a finished object's exterior can be specified without
specifying the deformation that produces it. (Section 4.8)

\textbf{Exposure condition.} The requirement
\(f(E_i) \subset \partial\Omega\), that a selected region ends up on the
finished object's apparent exterior boundary, within a stated tolerance.
(Section 4.2)

\textbf{Adjacency condition.} The requirement
\(A(E_i{:}e_a,\, E_j{:}e_b)\), that a specified edge of one selected
region meet a specified edge of another in the realized exterior.
(Section 4.3)

\textbf{Vertex-gluing condition.} The requirement
\(G(F_1, F_2, F_3) = V\), that designated corners of several local
frames become incident to one shared exterior vertex. (Section 4.4)

\textbf{Orientation condition.} The requirement that a realized
deformation carry a local frame's declared orientation into agreement
with its target orientation, up to tolerance. (Section 4.5)

\textbf{Relegation condition.} The requirement
\(f(I) \subset \Omega \cup H\), that the infoldable complement's image
lie within the object's interior or its hidden boundary layers \(H\),
rather than protruding as unaccounted-for exterior surface. (Section
4.6)

\textbf{Closure condition.} Any of a family of requirements governing
how complete, stable, or functionally container-like the finished object
must be; deliberately left as a family rather than a single formula,
since adequate closure depends on the construction's purpose. (Section
4.7)

\textbf{\(\text{INFOLD}(I \mid C)\).} The operator denoting constrained
search: deform the free region \(I\) through any permitted combination
of admissible operations such that the constraint set \(C\) is
satisfied, or is not thereby rendered unsatisfiable. Denotes a search
family, not a determinate procedure. (Section 5.1)

\textbf{Protected geometry.} Geometry corresponding directly to a
declared constraint, which a realization must not sacrifice in the
course of resolving anything else. (Section 5.4)

\textbf{Sacrificial geometry.} Geometry whose precise final shape and
crease history are not protected, and whose specificity, though not
whose material, may be surrendered in service of satisfying protected
constraints elsewhere. (Section 5.4)

\textbf{\(\text{LOCK}(C_k)\).} The operation of treating a currently
satisfied constraint as provisionally fixed, so that subsequent
operations are chosen or avoided specifically to preserve it. (Section
5.5)

\textbf{\(\text{RELAX}(C_k, \varepsilon)\).} The operation of permitting
a previously locked constraint to drift within a stated tolerance
\(\varepsilon\), undertaken to clear an obstruction elsewhere in the
construction. (Section 5.6)

\textbf{\(\text{REPAIR}(C_k)\).} The operation of reasserting a relaxed
constraint, ideally within its original tolerance, once the obstruction
that necessitated relaxation has been cleared. (Section 5.6)

\textbf{Sheet graph (\(G_0\)).} A graph whose nodes correspond to marked
regions and intervening portions of the flat, unfolded substrate, and
whose edges record direct connection across the original planar surface.
(Section 6.1)

\textbf{Boundary graph (\(G_\partial\)).} A graph whose nodes correspond
to regions visible on the finished exterior, and whose edges record
adjacency among them in the realized object. (Section 6.1)

\textbf{Surplus ratio (\(\sigma\)).} The quantity
\(\sigma = \big(\text{Area}(P) - \text{Area}(E_{\text{visible}})\big) / \text{Area}(E_{\text{visible}})\),
measuring how much of the substrate's area lies outside the declared
exterior, relative to the declared exterior's own area. (Section 6.4)

\textbf{Reachable set (\(\text{Reach}(P, h_t)\)).} The set of
configurations reachable from the substrate's current state, given a
specific deformation history \(h_t\) up to time \(t\); a property of the
pair of sheet and history, not of the sheet's current shape alone.
(Section 6.6)

\textbf{Boundary equivalence.} Two realizations satisfy the same
declared exposure, adjacency, orientation, vertex, and closure
conditions, regardless of anything else about them. The weakest and, for
this monograph's purposes, most important notion of equivalence among
realizations. (Section 6.8)

\textbf{Structural equivalence.} Two realizations are
boundary-equivalent and additionally share the same pattern of interior
layer contacts. (Section 6.8)

\textbf{Crease equivalence.} Two realizations are structurally
equivalent and additionally share the same crease network up to a stated
tolerance. (Section 6.8)

\textbf{Procedural equivalence.} Two realizations were produced by the
same sequence of operations in the same order; the strongest notion of
equivalence among realizations, and the one closest to the
reproducibility conventional sequential diagrams aim at. (Section 6.8)

\textbf{Local admissibility.} A move that does not irreparably violate a
protected constraint and does not damage the substrate's material
integrity, evaluated on its own terms without regard to its consequences
for later moves. (Section 9.1)

\textbf{Global reachability.} The continued existence, from a given
intermediate state, of at least one path to a configuration satisfying
every declared constraint. Distinct from, and not implied by, a sequence
of individually locally admissible moves. (Section 9.2)

\textbf{Hierarchy of protected distinctions.} The ranking material
continuity \(>\) selected-face identity \(>\) required incidence \(>\)
edge orientation \(>\) planarity \(>\) cosmetic regularity, used to
decide which relation to preserve when preserving one requires
sacrificing another. (Section 9.3)

\textbf{Ideal realizability.} A specification is ideally realizable if
some continuous deformation of an idealized, zero-thickness, infinitely
flexible surface satisfies it. (Section 9.9)

\textbf{Mechanical realizability.} A specification is mechanically
realizable if some physically plausible deformation of a sheet with
realistic thickness, stiffness, and plastic behavior satisfies it.
(Section 9.9)

\textbf{Practical realizability.} A specification is practically
realizable if a construction process can actually be found, within
reasonable effort, to locate a satisfying deformation, distinct from the
mere existence of one. (Section 9.9)

\hypertarget{appendix-c.-open-problems}{%
\chapter{Open Problems}\label{appendix-c.-open-problems}}

The following questions are left unresolved by this monograph. Each is
stated with a pointer to the chapter in which the relevant background
concept was introduced, but none is developed here beyond the statement
given.

\textbf{Realizability.} Beyond the single demonstrated case, what is the
actual scope of the conjecture stated in Section 9.7? Is there a
decision procedure, even in principle, for determining whether a given
boundary specification is realizable from a given substrate, or is the
question undecidable in general once crumpling is admitted as a
permitted operation? (Section 9.6, Section 9.7)

\textbf{Minimum surplus.} For a given target boundary specification,
what is the smallest surplus ratio \(\sigma\), as defined in Section
6.4, from which the specification remains realizable? Does this minimum
depend only on the specification's combinatorial structure, or also on
material properties in the sense of Section 9.9's distinction between
ideal and mechanical realizability?

\textbf{Internal complexity.} Given a realized object, can a meaningful
measure of its hidden interior's complexity be defined, and if so, does
exterior regularity reliably predict a high value of that measure, as
the tradeoff discussed in Section 6.5 would suggest, or are there
realizations that combine a regular exterior with a comparatively simple
interior?

\textbf{Diversity of interiors.} How many structurally distinct
interiors, in the sense of the equivalence hierarchy of Section 6.8, can
realize the same boundary specification in practice? The single sheet
examined in this monograph demonstrated at least two; the general shape
of this space, for a fixed specification, remains unexamined. (Section
10.3 proposes an experimental approach.)

\textbf{Material scaling.} How do the obstructions catalogued in Section
9.8, particularly those arising from thickness and plastic behavior,
scale with the size of the substrate and the size of the target object?
Does a sufficiently large sheet overcome obstructions that a smaller
sheet of the same material cannot?

\textbf{Verification.} Given a finished, unopened object, what can be
determined about whether it satisfies a claimed boundary specification
without destructively unfolding it, and how does the difficulty of this
verification compare to the difficulty of the original search that
produced the object? (Section 9.9, Section 10.4)

\textbf{Computational search.} Does a computational solver, searching
the space of admissible deformations directly, tend to discover the same
class of realizations a human constructor finds by hand, or does it
discover structurally different solutions unavailable to embodied,
trace-guided search of the kind described in Chapter 7? (Section 10.4)

\textbf{Robotic construction.} Does the constraint-level notation of
Chapter 4 carry sufficient information to guide a robotic manipulator
lacking the tactile feedback described in Chapter 7, or does successful
hand construction depend on forms of local judgment the notation does
not, and perhaps cannot, fully capture? (Section 10.4)

\textbf{The cortical analogy's limits.} Nothing in this monograph
establishes which, if any, of the mechanisms surveyed in Section 8.3 the
structural analogy of Chapter 8 actually bears on beyond the general
geometric point stated in Section 8.6. Whether the notion of a boundary
grammar, suitably adapted, could contribute any precise, testable
content to developmental biology, as opposed to serving only as an
illustrative analogy, is left entirely open.

\textbf{Expressive completeness.} Does the notation developed in Chapter
4, together with the operators of Chapter 5, suffice to state every
boundary specification of practical interest, or are there classes of
desired exterior relations that resist expression in this vocabulary and
require its extension?

\backmatter

\hypertarget{bibliography}{%
\chapter*{Bibliography}
\addcontentsline{toc}{chapter}{Bibliography}\label{bibliography}}

Cambou, Anne D., and Narayanan Menon. ``Three-Dimensional Structure of a
Sheet Crumpled into a Ball.'' \emph{Proceedings of the National Academy
of Sciences} 108, no. 36 (2011): 14741--14745.

Cerda, Enrique, and L. Mahadevan. ``Geometry and Physics of Wrinkling.''
\emph{Physical Review Letters} 90, no. 7 (2003): 074302.

Demaine, Erik D., and Joseph O'Rourke. \emph{Geometric Folding
Algorithms: Linkages, Origami, Polyhedra}. Cambridge: Cambridge
University Press, 2007.

Grassé, Pierre-Paul. ``La Reconstruction du Nid et les Coordinations
Interindividuelles chez \emph{Bellicositermes natalensis} et
\emph{Cubitermes} sp. La Théorie de la Stigmergie.'' \emph{Insectes
Sociaux} 6, no. 1 (1959): 41--80.

Harbin, Robert. \emph{Secrets of Origami: The Japanese Art of Paper
Folding}. London: Hodder and Stoughton, 1963.

Hull, Thomas C. \emph{Project Origami: Activities for Exploring
Mathematics}. 2nd ed.~Boca Raton: A K Peters/CRC Press, 2012.

Lang, Robert J. \emph{Origami Design Secrets: Mathematical Methods for
an Ancient Art}. 2nd ed.~Boca Raton: A K Peters/CRC Press, 2011.

Randlett, Samuel. \emph{The Art of Origami}. New York: E. P. Dutton,
1961.

Theraulaz, Guy, and Eric Bonabeau. ``A Brief History of Stigmergy.''
\emph{Artificial Life} 5, no. 2 (1999): 97--116.

Tallinen, Tuomas, Jun Young Chung, François Rousseau, Nadine Girard,
Julien Lefèvre, and L. Mahadevan. ``On the Growth and Form of Cortical
Convolutions.'' \emph{Nature Physics} 12, no. 6 (2016): 588--593.

Toro, Roberto, and Yves Burnod. ``A Morphogenetic Model for the
Development of Cortical Convolutions.'' \emph{Cerebral Cortex} 15, no.
12 (2005): 1900--1913.

Van Essen, David C. ``A Tension-Based Theory of Morphogenesis and
Compact Wiring in the Central Nervous System.'' \emph{Nature} 385, no.
6614 (1997): 313--318.

Witten, Thomas A. ``Stress Focusing in Elastic Sheets.'' \emph{Reviews
of Modern Physics} 79, no. 2 (2007): 643--675.

Zeman, Adam, Michaela Dewar, and Sergio Della Sala. ``Lives Without
Imagery --- Congenital Aphantasia.'' \emph{Cortex} 73 (2015): 378--380.

Zilles, Karl, Nicola Palomero-Gallagher, and Katrin Amunts.
``Development of Cortical Folding During Evolution and Ontogeny.''
\emph{Trends in Neurosciences} 36, no. 5 (2013): 275--284.

\end{document}

